\documentclass[8pt]{beamer}
\usetheme{Madrid}
\usecolortheme{default}
\usepackage{amsmath, amssymb, graphicx, array, multirow, booktabs, xcolor, tikz}
\usetikzlibrary{shapes,arrows,positioning,fit,backgrounds}

% Compact font settings
\setbeamerfont{title}{size=\Large}
\setbeamerfont{subtitle}{size=\small}
\setbeamerfont{author}{size=\scriptsize}
\setbeamerfont{date}{size=\scriptsize}
\setbeamerfont{frametitle}{size=\large}
\setbeamerfont{framesubtitle}{size=\normalsize}
\setbeamerfont{enumerate item}{size=\footnotesize}
\setbeamerfont{itemize item}{size=\footnotesize}
\setbeamerfont{block title}{size=\footnotesize}

% Compact spacing
\setlength{\itemsep}{0.08ex}
\setlength{\parskip}{0.08ex}
\setlength{\leftmargini}{0.3cm}
\setbeamersize{text margin left=3mm, text margin right=3mm}
\addtobeamertemplate{frame begin}{\vspace{-0.4cm}}{}

% Colors
\definecolor{myblue}{RGB}{0,0,255}
\definecolor{myred}{RGB}{255,0,0}
\definecolor{mygreen}{RGB}{0,128,0}
\definecolor{myorange}{RGB}{255,165,0}
\definecolor{mypurple}{RGB}{128,0,128}
\definecolor{mygray}{RGB}{128,128,128}
\definecolor{mygold}{RGB}{255,215,0}

\title{Supervised Learning - Model Performance Evaluation}
\subtitle{100 Questions: Continuous Metrics, Classification Metrics, ROC/AUC, Threshold Selection}
\author{GARP RAI Exam Preparation}
\date{}

\begin{document}
	
	\begin{frame}
		\titlepage
	\end{frame}
	
	% ============================================================
	% SECTION 1: INTRODUCTION TO MODEL EVALUATION - QUESTIONS 1-5
	% ============================================================
	
	\begin{frame}{Q1: Model Evaluation Purpose - Tutorial}
		\fontsize{6.5}{7.5}\selectfont
		\textbf{QUESTION: What is the primary purpose of model performance evaluation in supervised learning?}
		
		\vspace{0.1cm}
		\textbf{OPTIONS:}
		\begin{enumerate}
			\item To increase the complexity of the model
			\item To assess how well a model generalizes to new, unseen data and compare different models
			\item To add more features to the model
			\item To prepare the data for training
		\end{enumerate}
		
		\vspace{0.1cm}
		\textcolor{myblue}{\textbf{ANSWER: B}}
		
		\vspace{0.1cm}
		\textbf{DETAILED TUTORIAL:}
		\begin{itemize}
			\item \textbf{Primary purposes of model evaluation:}
			\begin{enumerate}
				\item \textbf{Generalization assessment:} How well does the model perform on new data?
				\item \textbf{Model comparison:} Which model is best for the specific problem?
				\item \textbf{Error quantification:} How wrong can we afford to be?
				\item \textbf{Risk management:} What are the consequences when models are wrong?
			\end{enumerate}
			\item \textbf{Why it matters in risk management:}
			\begin{itemize}
				\item Models drive critical decisions (credit, fraud, AML)
				\item Poor models have significant consequences
				\item Regulatory requirements demand validation
				\item Stakeholder confidence depends on reliable models
			\end{itemize}
		\end{itemize}
		
		\textcolor{myred}{\textbf{EXAM TRAP:}} Model evaluation is about \textcolor{myblue}{generalization performance}, not training performance!
	\end{frame}
	
	\begin{frame}{Q2: Continuous vs Classification - Tutorial}
		\fontsize{6.5}{7.5}\selectfont
		\textbf{QUESTION: Why is a distinction made between evaluation metrics for continuous vs discrete outputs?}
		
		\vspace{0.1cm}
		\textbf{OPTIONS:}
		\begin{enumerate}
			\item Because continuous outputs don't need evaluation
			\item Different types of outputs require different types of error measurement and interpretation
			\item Because classification models are always better
			\item Because continuous models are always more complex
		\end{enumerate}
		
		\vspace{0.1cm}
		\textcolor{myblue}{\textbf{ANSWER: B}}
		
		\vspace{0.1cm}
		\textbf{DETAILED TUTORIAL:}
		\begin{itemize}
			\item \textbf{Continuous outputs (regression):}
			\begin{itemize}
				\item Predict numbers: VaR, prices, revenues, losses
				\item Need measures of numerical accuracy
				\item Error magnitude matters
				\item Examples: MSE, RMSE, MAE, MAPE
			\end{itemize}
			\item \textbf{Discrete outputs (classification):}
			\begin{itemize}
				\item Predict categories: fraud vs legitimate, default vs paid
				\item Need measures of classification accuracy
				\item Type of error matters (false positives vs false negatives)
				\item Examples: Accuracy, Precision, Recall, F1, AUC
			\end{itemize}
		\end{itemize}
		
		\textcolor{myred}{\textbf{EXAM TRAP:}} \textcolor{myblue}{Continuous vs discrete} outputs require fundamentally different evaluation approaches!
	\end{frame}
	
	\begin{frame}{Q3: Poor Modeling Consequences - Tutorial}
		\fontsize{6.5}{7.5}\selectfont
		\textbf{QUESTION: What are the potential consequences of poor model performance in banking?}
		
		\vspace{0.1cm}
		\textbf{OPTIONS:}
		\begin{enumerate}
			\item Reduced capital requirements
			\item Capital add-ons, shareholder dilution, regulatory scrutiny, and reputational damage
			\item Decreased regulatory oversight
			\item Improved market confidence
		\end{enumerate}
		
		\vspace{0.1cm}
		\textcolor{myblue}{\textbf{ANSWER: B}}
		
		\vspace{0.1cm}
		\textbf{DETAILED TUTORIAL:}
		\begin{itemize}
			\item \textbf{Example from the chapter:}
			\begin{itemize}
				\item Greek bank estimated recovery rate at 65\%
				\item Actual recovery was 40\%
				\item 25-point overestimate forced €120 million capital add-on
				\item Triggered stock price slide
			\end{itemize}
			\item \textbf{Chain of consequences:}
			\begin{enumerate}
				\item \textbf{Financial impact:} Capital add-ons, shareholder dilution
				\item \textbf{Operational impact:} Deep-dive reviews, resource consumption
				\item \textbf{Strategic impact:} Revealed systemic weaknesses
				\item \textbf{Reputational impact:} Regulator and market confidence erosion
			\end{enumerate}
		\end{itemize}
		
		\textcolor{myred}{\textbf{EXAM TRAP:}} Poor modeling has \textcolor{myblue}{real financial and reputational consequences}!
	\end{frame}
	
	\begin{frame}{Q4: Test Sample Importance - Tutorial}
		\fontsize{6.5}{7.5}\selectfont
		\textbf{QUESTION: Why is it important to compute performance metrics on a test sample rather than the training sample?}
		
		\vspace{0.1cm}
		\textbf{OPTIONS:}
		\begin{enumerate}
			\item Because the test sample is larger
			\item Models often overfit to training data, so training performance is not a true test of generalization
			\item Because test samples are easier to compute
			\item Because training samples contain errors
		\end{enumerate}
		
		\vspace{0.1cm}
		\textcolor{myblue}{\textbf{ANSWER: B}}
		
		\vspace{0.1cm}
		\textbf{DETAILED TUTORIAL:}
		\begin{itemize}
			\item \textbf{Overfitting problem:}
			\begin{itemize}
				\item Models learn training data patterns
				\item Including noise and random fluctuations
				\item Performance on training data is optimistic
				\item Does not reflect true generalization
			\end{itemize}
			\item \textbf{Data split purposes:}
			\begin{itemize}
				\item \textbf{Training sample:} Estimate parameters
				\item \textbf{Validation sample:} Tune hyperparameters
				\item \textbf{Test sample:} Evaluate final performance
			\end{itemize}
		\end{itemize}
		
		\textcolor{myred}{\textbf{EXAM TRAP:}} Always evaluate on \textcolor{myblue}{test sample} for true generalization performance!
	\end{frame}
	
	\begin{frame}{Q5: Training vs Test Gap - Tutorial}
		\fontsize{6.5}{7.5}\selectfont
		\textbf{QUESTION: What does a significant gap between training and test performance indicate?}
		
		\vspace{0.1cm}
		\textbf{OPTIONS:}
		\begin{enumerate}
			\item The model is performing perfectly
			\item The model is likely overfitting to the training data
			\item The model is underfitting
			\item The data is perfectly clean
		\end{enumerate}
		
		\vspace{0.1cm}
		\textcolor{myblue}{\textbf{ANSWER: B}}
		
		\vspace{0.1cm}
		\textbf{DETAILED TUTORIAL:}
		\begin{itemize}
			\item \textbf{Performance patterns:}
			\begin{itemize}
				\item \textbf{No gap:} Training equal to Test (good generalization)
				\item \textbf{Small gap:} Slight overfitting (acceptable)
				\item \textbf{Large gap:} Significant overfitting (problematic)
			\end{itemize}
			\item \textbf{Why gap occurs:}
			\begin{itemize}
				\item Model memorizes training patterns
				\item Learns noise instead of signal
				\item Too complex for the data
				\item Insufficient regularization
			\end{itemize}
		\end{itemize}
		
		\textcolor{myred}{\textbf{EXAM TRAP:}} Large train-test performance gap = \textcolor{myblue}{overfitting}!
	\end{frame}
	
	% ============================================================
	% SECTION 2: CONTINUOUS OUTPUT METRICS - QUESTIONS 6-25
	% ============================================================
	
	\begin{frame}{Q6: MSE Definition - Tutorial}
		\fontsize{6.5}{7.5}\selectfont
		\textbf{QUESTION: What is Mean Squared Error (MSE) and what does it measure?}
		
		\vspace{0.1cm}
		\textbf{OPTIONS:}
		\begin{enumerate}
			\item Average of absolute differences between predictions and actual values
			\item Average of squared differences between predicted and actual values
			\item Percentage of correctly classified instances
			\item Average of squared differences divided by actual values
		\end{enumerate}
		
		\vspace{0.1cm}
		\textcolor{myblue}{\textbf{ANSWER: B}}
		
		\vspace{0.1cm}
		\textbf{DETAILED TUTORIAL:}
		\begin{itemize}
			\item \textbf{Formula:}
			\begin{itemize}
				\item $MSE = \frac{1}{N}\sum_{i=1}^N (y_i - \hat{y}_i)^2$
				\item $y_i$: true value
				\item $\hat{y}_i$: predicted value
				\item $N$: number of observations
			\end{itemize}
			\item \textbf{What it measures:}
			\begin{itemize}
				\item Average squared deviation
				\item Penalizes large errors more heavily
				\item Sensitive to outliers
			\end{itemize}
			\item \textbf{Key characteristics:}
			\begin{itemize}
				\item Always non-negative
				\item Lower is better
				\item Units are squared (e.g., dollars squared)
				\item Heavily penalizes large deviations
			\end{itemize}
		\end{itemize}
		
		\textcolor{myred}{\textbf{EXAM TRAP:}} MSE uses \textcolor{myblue}{squared errors} - heavily penalizes large deviations!
	\end{frame}
	
	\begin{frame}{Q7: MSE Managerial Question - Tutorial}
		\fontsize{6.5}{7.5}\selectfont
		\textbf{QUESTION: What managerial question does MSE help answer?}
		
		\vspace{0.1cm}
		\textbf{OPTIONS:}
		\begin{enumerate}
			\item "What does the typical error feel like in our primary unit of currency?"
			\item "How do we penalize big errors in our model?"
			\item "On average, by how much are our predictions off, regardless of direction?"
			\item "In percentage terms, what is the typical error relative to the actual value?"
		\end{enumerate}
		
		\vspace{0.1cm}
		\textcolor{myblue}{\textbf{ANSWER: B}}
		
		\vspace{0.1cm}
		\textbf{DETAILED TUTORIAL:}
		\begin{itemize}
			\item \textbf{Managerial focus:}
			\begin{itemize}
				\item Highlights tail risk
				\item Emphasis on squared errors heavily penalizes large deviations
				\item Sensitive to costly outliers
				\item Important for risk management
			\end{itemize}
			\item \textbf{Why this matters:}
			\begin{itemize}
				\item Large errors have significant consequences
				\item Risk managers care about tail events
				\item Regulatory focus on extreme outcomes
			\end{itemize}
		\end{itemize}
		
		\textcolor{myred}{\textbf{EXAM TRAP:}} MSE is for \textcolor{myblue}{penalizing big errors} - risk managers love it!
	\end{frame}
	
	\begin{frame}{Q8: RMSE Definition - Tutorial}
		\fontsize{6.5}{7.5}\selectfont
		\textbf{QUESTION: What is Root Mean Square Error (RMSE) and how does it differ from MSE?}
		
		\vspace{0.1cm}
		\textbf{OPTIONS:}
		\begin{enumerate}
			\item RMSE is the average of absolute errors
			\item Square root of MSE; scales with units of data, easier to interpret
			\item RMSE is always larger than MSE
			\item RMSE is the percentage of errors
		\end{enumerate}
		
		\vspace{0.1cm}
		\textcolor{myblue}{\textbf{ANSWER: B}}
		
		\vspace{0.1cm}
		\textbf{DETAILED TUTORIAL:}
		\begin{itemize}
			\item \textbf{Formula:}
			\begin{itemize}
				\item $RMSE = \sqrt{MSE} = \sqrt{\frac{1}{N}\sum_{i=1}^N (y_i - \hat{y}_i)^2}$
			\end{itemize}
			\item \textbf{Key characteristics:}
			\begin{itemize}
				\item Same units as target variable (e.g., dollars)
				\item More interpretable than MSE
				\item Gives sense of average error magnitude
				\item Still sensitive to outliers
			\end{itemize}
			\item \textbf{Comparison to MSE:}
			\begin{itemize}
				\item MSE: dollars squared (hard to interpret)
				\item RMSE: dollars (easy to interpret)
				\item Both penalize large errors
			\end{itemize}
		\end{itemize}
		
		\textcolor{myred}{\textbf{EXAM TRAP:}} RMSE is \textcolor{myblue}{easier to interpret} than MSE because it's in the same units!
	\end{frame}
	
	\begin{frame}{Q9: RMSE Managerial Question - Tutorial}
		\fontsize{6.5}{7.5}\selectfont
		\textbf{QUESTION: What managerial question does RMSE help answer?}
		
		\vspace{0.1cm}
		\textbf{OPTIONS:}
		\begin{enumerate}
			\item "How do we penalize big errors in our model?"
			\item "What does the typical error feel like in our primary unit of currency?"
			\item "On average, by how much are our predictions off?"
			\item "In percentage terms, what is the typical error?"
		\end{enumerate}
		
		\vspace{0.1cm}
		\textcolor{myblue}{\textbf{ANSWER: B}}
		
		\vspace{0.1cm}
		\textbf{DETAILED TUTORIAL:}
		\begin{itemize}
			\item \textbf{Managerial focus:}
			\begin{itemize}
				\item Same units as target variable (e.g., dollars)
				\item Directly interpretable
				\item "On average, our predictions are off by RMSE dollars"
				\item Helps non-technical stakeholders understand
			\end{itemize}
			\item \textbf{Caveat:}
			\begin{itemize}
				\item Still sensitive to outliers
				\item Unusually large errors can inflate RMSE
				\item Need to investigate if high RMSE is systemic or from extreme events
			\end{itemize}
		\end{itemize}
		
		\textcolor{myred}{\textbf{EXAM TRAP:}} RMSE answers \textcolor{myblue}{"what does the typical error feel like?"}
	\end{frame}
	
	\begin{frame}{Q10: MAE Definition - Tutorial}
		\fontsize{6.5}{7.5}\selectfont
		\textbf{QUESTION: What is Mean Absolute Error (MAE) and what is its key advantage?}
		
		\vspace{0.1cm}
		\textbf{OPTIONS:}
		\begin{enumerate}
			\item Average of squared errors; penalizes large errors
			\item Average of absolute differences; robust to outliers
			\item Average of percentage errors; good for comparisons
			\item Average of errors; can cancel out positive and negative
		\end{enumerate}
		
		\vspace{0.1cm}
		\textcolor{myblue}{\textbf{ANSWER: B}}
		
		\vspace{0.1cm}
		\textbf{DETAILED TUTORIAL:}
		\begin{itemize}
			\item \textbf{Formula:}
			\begin{itemize}
				\item $MAE = \frac{1}{N}\sum_{i=1}^N |y_i - \hat{y}_i|$
			\end{itemize}
			\item \textbf{Key characteristics:}
			\begin{itemize}
				\item Uses absolute differences (not squared)
				\item Robust to outliers
				\item Easy to communicate
				\item Straightforward average error magnitude
			\end{itemize}
			\item \textbf{Limitations:}
			\begin{itemize}
				\item Can mask direction of bias
				\item Doesn't penalize occasional large misses
				\item May underestimate critical risk events
			\end{itemize}
		\end{itemize}
		
		\textcolor{myred}{\textbf{EXAM TRAP:}} MAE is \textcolor{myblue}{robust to outliers} - unlike MSE and RMSE!
	\end{frame}
	
	\begin{frame}{Q11: MAE Managerial Question - Tutorial}
		\fontsize{6.5}{7.5}\selectfont
		\textbf{QUESTION: What managerial question does MAE help answer?}
		
		\vspace{0.1cm}
		\textbf{OPTIONS:}
		\begin{enumerate}
			\item "How do we penalize big errors?"
			\item "What does the typical error feel like?"
			\item "On average, by how much are our predictions off, regardless of direction?"
			\item "In percentage terms, what is the typical error?"
		\end{enumerate}
		
		\vspace{0.1cm}
		\textcolor{myblue}{\textbf{ANSWER: C}}
		
		\vspace{0.1cm}
		\textbf{DETAILED TUTORIAL:}
		\begin{itemize}
			\item \textbf{Managerial focus:}
			\begin{itemize}
				\item "On average, our predictions are off by MAE dollars"
				\item Ignores direction (over vs under prediction)
				\item Easy to communicate to non-technical audiences
			\end{itemize}
			\item \textbf{When to use:}
			\begin{itemize}
				\item When outliers are not critical
				\item When you want stable, interpretable measure
				\item When equal treatment of all errors is desired
			\end{itemize}
		\end{itemize}
		
		\textcolor{myred}{\textbf{EXAM TRAP:}} MAE answers \textcolor{myblue}{"by how much are we off, regardless of direction?"}
	\end{frame}
	
	\begin{frame}{Q12: MAPE Definition - Tutorial}
		\fontsize{6.5}{7.5}\selectfont
		\textbf{QUESTION: What is Mean Absolute Percentage Error (MAPE) and what is its key advantage?}
		
		\vspace{0.1cm}
		\textbf{OPTIONS:}
		\begin{enumerate}
			\item Average of squared percentage errors
			\item Average of absolute percentage errors; provides normalized percentage measure
			\item Average of percentage errors; can cancel out
			\item Average of absolute errors divided by predictions
		\end{enumerate}
		
		\vspace{0.1cm}
		\textcolor{myblue}{\textbf{ANSWER: B}}
		
		\vspace{0.1cm}
		\textbf{DETAILED TUTORIAL:}
		\begin{itemize}
			\item \textbf{Formula:}
			\begin{itemize}
				\item $MAPE = \frac{100}{N}\sum_{i=1}^N \left|\frac{y_i - \hat{y}_i}{y_i}\right|$
			\end{itemize}
			\item \textbf{Key characteristics:}
			\begin{itemize}
				\item Expressed as percentage
				\item Normalized - allows comparison across different scales
				\item Board-friendly (percentage nature)
				\item Consistent across portfolios
			\end{itemize}
			\item \textbf{Limitations:}
			\begin{itemize}
				\item Explodes when actual values are zero or near zero
				\item Can be misleading with small denominators
				\item May understate errors on large absolute value items
			\end{itemize}
		\end{itemize}
		
		\textcolor{myred}{\textbf{EXAM TRAP:}} MAPE is \textcolor{myblue}{percentage-based} - great for communication but has limitations!
	\end{frame}
	
	\begin{frame}{Q13: MAPE Managerial Question - Tutorial}
		\fontsize{6.5}{7.5}\selectfont
		\textbf{QUESTION: What managerial question does MAPE help answer?}
		
		\vspace{0.1cm}
		\textbf{OPTIONS:}
		\begin{enumerate}
			\item "How do we penalize big errors?"
			\item "What does the typical error feel like?"
			\item "On average, by how much are our predictions off?"
			\item "In percentage terms, what is the typical error relative to the actual value?"
		\end{enumerate}
		
		\vspace{0.1cm}
		\textcolor{myblue}{\textbf{ANSWER: D}}
		
		\vspace{0.1cm}
		\textbf{DETAILED TUTORIAL:}
		\begin{itemize}
			\item \textbf{Managerial focus:}
			\begin{itemize}
				\item "On average, our predictions are off by MAPE\%"
				\item Normalized - allows comparison across models/portfolios
				\item Perfect for audit committees and board presentations
				\item Common choice for auditors
			\end{itemize}
			\item \textbf{When to use:}
			\begin{itemize}
				\item Comparing models with different scales
				\item Communicating to non-technical stakeholders
				\item When percentage errors are more meaningful than absolute
			\end{itemize}
		\end{itemize}
		
		\textcolor{myred}{\textbf{EXAM TRAP:}} MAPE answers \textcolor{myblue}{"what's the error as a percentage?"}
	\end{frame}
	
	\begin{frame}{Q14: MSE Calculation Example - Tutorial}
		\fontsize{6.5}{7.5}\selectfont
		\textbf{QUESTION: Given a test sample of 10 observations, how is MSE calculated?}
		
		\vspace{0.1cm}
		\textbf{OPTIONS:}
		\begin{enumerate}
			\item Sum of absolute errors divided by N
			\item Sum of squared errors divided by N
			\item Square root of sum of squared errors divided by N
			\item Sum of errors divided by N
		\end{enumerate}
		
		\vspace{0.1cm}
		\textcolor{myblue}{\textbf{ANSWER: B}}
		
		\vspace{0.1cm}
		\textbf{DETAILED TUTORIAL:}
		\begin{itemize}
			\item \textbf{Step-by-step with example:}
			\begin{enumerate}
				\item Predict values using the model
				\item Calculate errors: $y_i - \hat{y}_i$
				\item Square each error: $(y_i - \hat{y}_i)^2$
				\item Sum all squared errors
				\item Divide by N (number of observations)
			\end{enumerate}
			\item \textbf{Example from chapter:}
			\begin{itemize}
				\item Salary prediction model: $\hat{y} = 16.88 + 1.06 \times \text{experience}$
				\item Sum of squared errors: 431.02
				\item N = 10
				\item MSE = 431.02 / 10 = 43.1
				\item RMSE = $\sqrt{43.1} = 6.57$
			\end{itemize}
		\end{itemize}
		
		\textcolor{myred}{\textbf{EXAM TRAP:}} MSE = \textcolor{myblue}{sum of squared errors divided by N}!
	\end{frame}
	
	\begin{frame}{Q15: RMSE Interpretation - Tutorial}
		\fontsize{6.5}{7.5}\selectfont
		\textbf{QUESTION: In the salary prediction example with RMSE = 6.57, what does this mean?}
		
		\vspace{0.1cm}
		\textbf{OPTIONS:}
		\begin{enumerate}
			\item The model is off by exactly 6.57 for every prediction
			\item On average, the model's predictions are off by about 6.57 units
			\item The model is 6.57\% accurate
			\item The model has 6.57\% error rate
		\end{enumerate}
		
		\vspace{0.1cm}
		\textcolor{myblue}{\textbf{ANSWER: B}}
		
		\vspace{0.1cm}
		\textbf{DETAILED TUTORIAL:}
		\begin{itemize}
			\item \textbf{Interpretation:}
			\begin{itemize}
				\item "On average, predictions are off by about 6.57 thousand dollars"
				\item RMSE is in same units as the target (salary in \$000s)
				\item Gives sense of typical error magnitude
			\end{itemize}
			\item \textbf{Important nuances:}
			\begin{itemize}
				\item This is an average - some errors are larger, some smaller
				\item Squaring means larger errors contribute more to the average
				\item RMSE eq 6.57 means typical error is around 6.57 units
			\end{itemize}
		\end{itemize}
		
		\textcolor{myred}{\textbf{EXAM TRAP:}} RMSE is in \textcolor{myblue}{same units as target} - easier to understand!
	\end{frame}
	
	\begin{frame}{Q16: MAE Example Calculation - Tutorial}
		\fontsize{6.5}{7.5}\selectfont
		\textbf{QUESTION: In the salary prediction example, MAE = 5.01. What does this mean and how is it calculated?}
		
		\vspace{0.1cm}
		\textbf{OPTIONS:}
		\begin{enumerate}
			\item Average of squared errors = 5.01
			\item Average of absolute errors = 5.01 (sum of absolute errors divided by N)
			\item Average of percentage errors = 5.01\%
			\item Square root of MSE = 5.01
		\end{enumerate}
		
		\vspace{0.1cm}
		\textcolor{myblue}{\textbf{ANSWER: B}}
		
		\vspace{0.1cm}
		\textbf{DETAILED TUTORIAL:}
		\begin{itemize}
			\item \textbf{Calculation:}
			\begin{itemize}
				\item Sum of absolute errors = 50.11 (from Table 8.3)
				\item N = 10
				\item MAE = 50.11 / 10 = 5.01
			\end{itemize}
			\item \textbf{Interpretation:}
			\begin{itemize}
				\item "On average, predictions are off by about 5.01 thousand dollars"
				\item Similar to RMSE but less influenced by large errors
				\item More stable measure of typical error
			\end{itemize}
		\end{itemize}
		
		\textcolor{myred}{\textbf{EXAM TRAP:}} MAE = \textcolor{myblue}{average of absolute errors} - sum divided by N!
	\end{frame}
	
	\begin{frame}{Q17: MAPE Example Calculation - Tutorial}
		\fontsize{6.5}{7.5}\selectfont
		\textbf{QUESTION: In the salary example, MAPE = 16\%. What does this mean?}
		
		\vspace{0.1cm}
		\textbf{OPTIONS:}
		\begin{enumerate}
			\item The model is 16\% accurate
			\item The average forecast error is 16\% of actual salary
			\item The model makes errors 16\% of the time
			\item The model has 16\% bias
		\end{enumerate}
		
		\vspace{0.1cm}
		\textcolor{myblue}{\textbf{ANSWER: B}}
		
		\vspace{0.1cm}
		\textbf{DETAILED TUTORIAL:}
		\begin{itemize}
			\item \textbf{Calculation:}
			\begin{itemize}
				\item Sum of $|(y_i - \hat{y}_i)/y_i|$ = 1.62
				\item N = 10
				\item MAPE = 1.62 / 10 × 100 = 16.2\%
			\end{itemize}
			\item \textbf{Interpretation:}
			\begin{itemize}
				\item "On average, the forecast error is 16\% of the actual salary"
				\item Most intuitive of all measures
				\item Easy to communicate to non-technical stakeholders
			\end{itemize}
		\end{itemize}
		
		\textcolor{myred}{\textbf{EXAM TRAP:}} MAPE = \textcolor{myblue}{percentage of actual values} - most intuitive for communication!
	\end{frame}
	
	\begin{frame}{Q18: MSE vs RMSE vs MAE vs MAPE - Tutorial}
		\fontsize{6.5}{7.5}\selectfont
		\textbf{QUESTION: Which metric would a risk manager choose to capture tail risk?}
		
		\vspace{0.1cm}
		\textbf{OPTIONS:}
		\begin{enumerate}
			\item MAE
			\item MSE or RMSE
			\item MAPE
			\item Accuracy
		\end{enumerate}
		
		\vspace{0.1cm}
		\textcolor{myblue}{\textbf{ANSWER: B}}
		
		\vspace{0.1cm}
		\textbf{DETAILED TUTORIAL:}
		\begin{itemize}
			\item \textbf{Why MSE/RMSE capture tail risk:}
			\begin{itemize}
				\item Squaring errors heavily penalizes large deviations
				\item Large errors contribute disproportionately
				\item Sensitive to outliers (tail events)
				\item Regulatory preference (ECB recommends)
			\end{itemize}
			\item \textbf{Comparison:}
			\begin{itemize}
				\item MAE: Robust, but ignores tail risk
				\item MAPE: Good for communication, but has limitations
				\item MSE/RMSE: Best for capturing extreme errors
			\end{itemize}
		\end{itemize}
		
		\textcolor{myred}{\textbf{EXAM TRAP:}} ECB recommends metrics that \textcolor{myblue}{capture tail risk} - MSE and RMSE!
	\end{frame}
	
	\begin{frame}{Q19: Metric Selection - Tutorial}
		\fontsize{6.5}{7.5}\selectfont
		\textbf{QUESTION: When should a risk manager use MAPE instead of MSE or RMSE?}
		
		\vspace{0.1cm}
		\textbf{OPTIONS:}
		\begin{enumerate}
			\item When comparing models with very different scales
			\item When tail risk is the primary concern
			\item When errors are always very small
			\item When the model is overfitting
		\end{enumerate}
		
		\vspace{0.1cm}
		\textcolor{myblue}{\textbf{ANSWER: A}}
		
		\vspace{0.1cm}
		\textbf{DETAILED TUTORIAL:}
		\begin{itemize}
			\item \textbf{When MAPE is appropriate:}
			\begin{itemize}
				\item Comparing models across different portfolios
				\item Comparing a €100k loan vs €10m loan
				\item Communicating to audit committees
				\item When percentage errors are more meaningful
			\end{itemize}
			\item \textbf{When MAPE is NOT appropriate:}
			\begin{itemize}
				\item When actual values are near zero
				\item When tail risk is the primary concern
				\item When absolute errors matter more than percentage
			\end{itemize}
		\end{itemize}
		
		\textcolor{myred}{\textbf{EXAM TRAP:}} MAPE is best for \textcolor{myblue}{comparing across different scales}!
	\end{frame}
	
	\begin{frame}{Q20: Error Distribution Visualization - Tutorial}
		\fontsize{6.5}{7.5}\selectfont
		\textbf{QUESTION: Why should risk managers look beyond average metrics and examine error distributions?}
		
		\vspace{0.1cm}
		\textbf{OPTIONS:}
		\begin{enumerate}
			\item Average metrics are always sufficient
			\item The distribution can reveal patterns or biases that averages might hide
			\item Error distributions are easier to calculate
			\item Averages are not used in risk management
		\end{enumerate}
		
		\vspace{0.1cm}
		\textcolor{myblue}{\textbf{ANSWER: B}}
		
		\vspace{0.1cm}
		\textbf{DETAILED TUTORIAL:}
		\begin{itemize}
			\item \textbf{Key questions to ask:}
			\begin{itemize}
				\item "Are errors generally small with a few large ones?"
				\item "Are errors randomly distributed or systematically biased?"
				\item "Are there segments where errors are consistently larger?"
			\end{itemize}
			\item \textbf{Why this matters:}
			\begin{itemize}
				\item Consistent bias > random errors
				\item Segments with larger errors need attention
				\item Visualizations (histograms, box plots) provide richer context
			\end{itemize}
		\end{itemize}
		
		\textcolor{myred}{\textbf{EXAM TRAP:}} Error distributions reveal \textcolor{myblue}{patterns and biases} that averages hide!
	\end{frame}
	
	\begin{frame}{Q21: Model Comparison Example - Tutorial}
		\fontsize{6.5}{7.5}\selectfont
		\textbf{QUESTION: In the house price prediction example, how do we determine which model is better?}
		
		\vspace{0.1cm}
		\textbf{OPTIONS:}
		\begin{enumerate}
			\item Always choose the more complex model
			\item Compare performance metrics on test sample - lower values indicate better performance
			\item Compare performance metrics on training sample
			\item Always choose the simpler model
		\end{enumerate}
		
		\vspace{0.1cm}
		\textcolor{myblue}{\textbf{ANSWER: B}}
		
		\vspace{0.1cm}
		\textbf{DETAILED TUTORIAL:}
		\begin{itemize}
			\item \textbf{Example results (Table 8.4):}
			\begin{itemize}
				\item Model A (linear regression): MSE=94.85, MAE=6.56, MAPE=18.25
				\item Model B (decision tree): MSE=77.35*, MAE=5.41*, MAPE=14.45*
			\end{itemize}
			\item \textbf{Conclusion:}
			\begin{itemize}
				\item Model B has all lower values
				\item Decision tree is better according to all metrics
				\item Asterisk denotes best performer
			\end{itemize}
		\end{itemize}
		
		\textcolor{myred}{\textbf{EXAM TRAP:}} Lower values = \textcolor{myblue}{better performance} for all these metrics!
	\end{frame}
	
	\begin{frame}{Q22: Different Metrics Different Rankings - Tutorial}
		\fontsize{6.5}{7.5}\selectfont
		\textbf{QUESTION: Can different performance metrics yield different model rankings?}
		
		\vspace{0.1cm}
		\textbf{OPTIONS:}
		\begin{enumerate}
			\item No, all metrics always give the same ranking
			\item Yes, different metrics can yield different rankings; choice depends on the problem
			\item Only MSE and RMSE give different rankings
			\item Only MAE and MAPE give different rankings
		\end{enumerate}
		
		\vspace{0.1cm}
		\textcolor{myblue}{\textbf{ANSWER: B}}
		
		\vspace{0.1cm}
		\textbf{DETAILED TUTORIAL:}
		\begin{itemize}
			\item \textbf{Why rankings can differ:}
			\begin{itemize}
				\item MSE/RMSE: Penalize large errors
				\item MAE: Treats all errors equally
				\item MAPE: Emphasizes relative errors
				\item Different metrics = different priorities
			\end{itemize}
			\item \textbf{Implication:}
			\begin{itemize}
				\item Choose metric based on business needs
				\item Regulatory requirements may dictate choice
				\item Historical organizational choices matter
			\end{itemize}
		\end{itemize}
		
		\textcolor{myred}{\textbf{EXAM TRAP:}} Different metrics = \textcolor{myblue}{different rankings} - choose based on problem!
	\end{frame}
	
	\begin{frame}{Q23: Unnormalized Metrics Limitation - Tutorial}
		\fontsize{6.5}{7.5}\selectfont
		\textbf{QUESTION: What is a limitation of unnormalized metrics like RMSE and MAE?}
		
		\vspace{0.1cm}
		\textbf{OPTIONS:}
		\begin{enumerate}
			\item They are too difficult to calculate
			\item They treat errors symmetrically regardless of the scale of the target
			\item They always prefer complex models
			\item They cannot be used for risk management
		\end{enumerate}
		
		\vspace{0.1cm}
		\textcolor{myblue}{\textbf{ANSWER: B}}
		
		\vspace{0.1cm}
		\textbf{DETAILED TUTORIAL:}
		\begin{itemize}
			\item \textbf{The problem:}
			\begin{itemize}
				\item \$148,000 vs \$150,000 = \$2,000 error (1%)
				\item \$10,000 vs \$12,000 = \$2,000 error (>10%)
				\item Both contribute same to RMSE/MAE
				\item But the second error is much more significant
			\end{itemize}
			\item \textbf{Solution:}
			\begin{itemize}
				\item Use MAPE for relative comparison
				\item Understand the context of errors
				\item Consider business impact, not just statistical measures
			\end{itemize}
		\end{itemize}
		
		\textcolor{myred}{\textbf{EXAM TRAP:}} Unnormalized metrics treat \textcolor{myblue}{all errors equally} - MAPE handles this issue!
	\end{frame}
	
	\begin{frame}{Q24: Regulatory Metric Preferences - Tutorial}
		\fontsize{6.5}{7.5}\selectfont
		\textbf{QUESTION: What is the ECB's recommendation regarding performance metrics?}
		
		\vspace{0.1cm}
		\textbf{OPTIONS:}
		\begin{enumerate}
			\item Only use MAPE for all models
			\item Use metrics that capture tail risk, nudging firms to report RMSE alongside MAE
			\item Always use MAE because it's most robust
			\item Never use MSE because it's too complex
		\end{enumerate}
		
		\vspace{0.1cm}
		\textcolor{myblue}{\textbf{ANSWER: B}}
		
		\vspace{0.1cm}
		\textbf{DETAILED TUTORIAL:}
		\begin{itemize}
			\item \textbf{ECB Guide to Internal Models:}
			\begin{itemize}
				\item Recommends metrics that capture tail risk
				\item Nudges firms to report RMSE alongside MAE
				\item Ensures large, infrequent losses are not overlooked
			\end{itemize}
			\item \textbf{Why this matters:}
			\begin{itemize}
				\item MAE alone might mask critical risks
				\item RMSE captures the impact of extreme events
				\item Both together provide complete picture
			\end{itemize}
		\end{itemize}
		
		\textcolor{myred}{\textbf{EXAM TRAP:}} ECB recommends \textcolor{myblue}{RMSE alongside MAE} for complete risk picture!
	\end{frame}
	
	\begin{frame}{Q25: Auditor Preferences - Tutorial}
		\fontsize{6.5}{7.5}\selectfont
		\textbf{QUESTION: Why do auditors often gravitate to MAPE?}
		
		\vspace{0.1cm}
		\textbf{OPTIONS:}
		\begin{enumerate}
			\item Because MAPE is most accurate
			\item For its pure percentage narrative when discussing model accuracy with audit committees
			\item Because MAPE captures tail risk best
			\item Because MAPE is easiest to calculate
		\end{enumerate}
		
		\vspace{0.1cm}
		\textcolor{myblue}{\textbf{ANSWER: B}}
		
		\vspace{0.1cm}
		\textbf{DETAILED TUTORIAL:}
		\begin{itemize}
			\item \textbf{Auditor perspective:}
			\begin{itemize}
				\item MAPE provides standardized percentage measure
				\item Easy to understand for non-technical stakeholders
				\item Allows comparison across different models/portfolios
				\item Communicates effectively with audit committees
			\end{itemize}
			\item \textbf{Why this matters:}
			\begin{itemize}
				\item Different stakeholders prefer different metrics
				\item Regulators vs Auditors have different needs
				\item Risk managers must satisfy multiple audiences
			\end{itemize}
		\end{itemize}
		
		\textcolor{myred}{\textbf{EXAM TRAP:}} Auditors prefer MAPE for its \textcolor{myblue}{percentage narrative}! 
	\end{frame}
	
	% ============================================================
	% SECTION 3: CONFUSION MATRIX - QUESTIONS 26-45
	% ============================================================
	
	\begin{frame}{Q26: Confusion Matrix Definition - Tutorial}
		\fontsize{6.5}{7.5}\selectfont
		\textbf{QUESTION: What is a confusion matrix in classification model evaluation?}
		
		\vspace{0.1cm}
		\textbf{OPTIONS:}
		\begin{enumerate}
			\item A matrix showing model parameters
			\item A cross-tabulation of predicted classifications against actual outcomes
			\item A matrix showing feature correlations
			\item A matrix showing error rates
		\end{enumerate}
		
		\vspace{0.1cm}
		\textcolor{myblue}{\textbf{ANSWER: B}}
		
		\vspace{0.1cm}
		\textbf{DETAILED TUTORIAL:}
		\begin{itemize}
			\item \textbf{Structure:}
			\begin{tabular}{|c|c|c|}
				\hline
				& Predicted Negative & Predicted Positive \\
				\hline
				Actual Negative & True Negative (TN) & False Positive (FP) \\
				\hline
				Actual Positive & False Negative (FN) & True Positive (TP) \\
				\hline
			\end{tabular}
			\item \textbf{Why it's important:}
			\begin{itemize}
				\item Foundation for all classification metrics
				\item Shows all four types of outcomes
				\item Enables calculation of Precision, Recall, Accuracy, F1
			\end{itemize}
		\end{itemize}
		
		\textcolor{myred}{\textbf{EXAM TRAP:}} Confusion matrix = \textcolor{myblue}{predicted vs actual} cross-tabulation!
	\end{frame}
	
	\begin{frame}{Q27: True Positive - Tutorial}
		\fontsize{6.5}{7.5}\selectfont
		\textbf{QUESTION: What is a True Positive (TP) in a confusion matrix?}
		
		\vspace{0.1cm}
		\textbf{OPTIONS:}
		\begin{enumerate}
			\item Correctly identified a negative case
			\item Correctly identified a positive case
			\item Incorrectly identified a negative case as positive
			\item Incorrectly identified a positive case as negative
		\end{enumerate}
		
		\vspace{0.1cm}
		\textcolor{myblue}{\textbf{ANSWER: B}}
		
		\vspace{0.1cm}
		\textbf{DETAILED TUTORIAL:}
		\begin{itemize}
			\item \textbf{Definition:} Model correctly identified a positive case
			\item \textbf{Examples:}
			\begin{itemize}
				\item Correctly flagged a fraudulent transaction
				\item Correctly predicted that a loan would default
				\item Correctly identified a dividend-paying firm
			\end{itemize}
			\item \textbf{Desired outcome:}
			\begin{itemize}
				\item This is what we want the model to do
				\item More TP = better model performance
			\end{itemize}
		\end{itemize}
		
		\textcolor{myred}{\textbf{EXAM TRAP:}} True Positive = \textcolor{myblue}{correctly identified positive}!
	\end{frame}
	
	\begin{frame}{Q28: True Negative - Tutorial}
		\fontsize{6.5}{7.5}\selectfont
		\textbf{QUESTION: What is a True Negative (TN) in a confusion matrix?}
		
		\vspace{0.1cm}
		\textbf{OPTIONS:}
		\begin{enumerate}
			\item Correctly identified a negative case
			\item Correctly identified a positive case
			\item Incorrectly identified a negative case as positive
			\item Incorrectly identified a positive case as negative
		\end{enumerate}
		
		\vspace{0.1cm}
		\textcolor{myblue}{\textbf{ANSWER: A}}
		
		\vspace{0.1cm}
		\textbf{DETAILED TUTORIAL:}
		\begin{itemize}
			\item \textbf{Definition:} Model correctly identified a negative case
			\item \textbf{Examples:}
			\begin{itemize}
				\item Correctly identified a legitimate transaction
				\item Correctly predicted that a loan would not default
				\item Correctly identified a non-dividend paying firm
			\end{itemize}
			\item \textbf{Desired outcome:}
			\begin{itemize}
				\item This is what we want the model to do
				\item More TN = better model performance
			\end{itemize}
		\end{itemize}
		
		\textcolor{myred}{\textbf{EXAM TRAP:}} True Negative = \textcolor{myblue}{correctly identified negative}!
	\end{frame}
	
	\begin{frame}{Q29: False Positive - Tutorial}
		\fontsize{6.5}{7.5}\selectfont
		\textbf{QUESTION: What is a False Positive (FP) in a confusion matrix?}
		
		\vspace{0.1cm}
		\textbf{OPTIONS:}
		\begin{enumerate}
			\item Correctly identified a negative case
			\item Correctly identified a positive case
			\item Incorrectly classified a negative case as positive (Type I Error)
			\item Incorrectly classified a positive case as negative (Type II Error)
		\end{enumerate}
		
		\vspace{0.1cm}
		\textcolor{myblue}{\textbf{ANSWER: C}}
		
		\vspace{0.1cm}
		\textbf{DETAILED TUTORIAL:}
		\begin{itemize}
			\item \textbf{Definition:} Model incorrectly classified a negative case as positive
			\item \textbf{Also known as:} Type I Error
			\item \textbf{Examples:}
			\begin{itemize}
				\item Flagged a legitimate transaction as fraudulent
				\item Predicted that a good loan would default
				\item False alarm
			\end{itemize}
			\item \textbf{Cost perspective:}
			\begin{itemize}
				\item Operational cost to investigate
				\item Customer friction
				\item Lost revenue from declined good customers
			\end{itemize}
		\end{itemize}
		
		\textcolor{myred}{\textbf{EXAM TRAP:}} False Positive = \textcolor{myblue}{Type I Error} - incorrectly saying "yes"!
	\end{frame}
	
	\begin{frame}{Q30: False Negative - Tutorial}
		\fontsize{6.5}{7.5}\selectfont
		\textbf{QUESTION: What is a False Negative (FN) in a confusion matrix?}
		
		\vspace{0.1cm}
		\textbf{OPTIONS:}
		\begin{enumerate}
			\item Correctly identified a negative case
			\item Correctly identified a positive case
			\item Incorrectly classified a negative case as positive
			\item Incorrectly classified a positive case as negative (Type II Error)
		\end{enumerate}
		
		\vspace{0.1cm}
		\textcolor{myblue}{\textbf{ANSWER: D}}
		
		\vspace{0.1cm}
		\textbf{DETAILED TUTORIAL:}
		\begin{itemize}
			\item \textbf{Definition:} Model incorrectly classified a positive case as negative
			\item \textbf{Also known as:} Type II Error
			\item \textbf{Examples:}
			\begin{itemize}
				\item Missed a fraudulent transaction (labeled legitimate)
				\item Failed to predict that a loan would default
				\item Missed a critical compliance breach
			\end{itemize}
			\item \textbf{Cost perspective:}
			\begin{itemize}
				\item Often the most costly type of error
				\item Financial losses from undetected fraud
				\item Capital impact of missed defaults
				\item Regulatory sanctions
			\end{itemize}
		\end{itemize}
		
		\textcolor{myred}{\textbf{EXAM TRAP:}} False Negative = \textcolor{myblue}{Type II Error} - incorrectly saying "no"!
	\end{frame}
	
	\begin{frame}{Q31: Type I vs Type II Errors - Tutorial}
		\fontsize{6.5}{7.5}\selectfont
		\textbf{QUESTION: Which type of error is typically more costly in credit risk management?}
		
		\vspace{0.1cm}
		\textbf{OPTIONS:}
		\begin{enumerate}
			\item Type I Error (False Positive)
			\item Type II Error (False Negative)
			\item Both are equally costly
			\item Neither is costly
		\end{enumerate}
		
		\vspace{0.1cm}
		\textcolor{myblue}{\textbf{ANSWER: B}}
		
		\vspace{0.1cm}
		\textbf{DETAILED TUTORIAL:}
		\begin{itemize}
			\item \textbf{Why Type II Errors are more costly:}
			\begin{itemize}
				\item Misclassifying a borrower as solvent who defaults
				\item Results in actual loan losses
				\item Regulatory capital implications
				\item Direct financial impact
			\end{itemize}
			\item \textbf{Type I Error cost:}
			\begin{itemize}
				\item Lost profit from not making a loan
				\item Operational cost of investigation
				\item Customer friction
				\item Usually less severe than Type II
			\end{itemize}
			\item \textbf{General principle:}
			\begin{itemize}
				\item The cost of errors is typically not symmetric
				\item Choose threshold based on relative costs
			\end{itemize}
		\end{itemize}
		
		\textcolor{myred}{\textbf{EXAM TRAP:}} Type II Errors (False Negatives) are \textcolor{myblue}{usually more costly} in credit risk!
	\end{frame}
	
	\begin{frame}{Q32: Accuracy Definition - Tutorial}
		\fontsize{6.5}{7.5}\selectfont
		\textbf{QUESTION: What is accuracy in classification model evaluation?}
		
		\vspace{0.1cm}
		\textbf{OPTIONS:}
		\begin{enumerate}
			\item Percentage of correctly classified positive cases
			\item Percentage of correctly classified instances overall
			\item Percentage of negative cases correctly classified
			\item Average of precision and recall
		\end{enumerate}
		
		\vspace{0.1cm}
		\textcolor{myblue}{\textbf{ANSWER: B}}
		
		\vspace{0.1cm}
		\textbf{DETAILED TUTORIAL:}
		\begin{itemize}
			\item \textbf{Formula:}
			\begin{itemize}
				\item $Accuracy = \frac{TP + TN}{TP + TN + FP + FN}$
			\end{itemize}
			\item \textbf{Interpretation:}
			\begin{itemize}
				\item "Percentage of all predictions that were correct"
				\item Most intuitive classification metric
				\item Easy to communicate
			\end{itemize}
			\item \textbf{Limitations:}
			\begin{itemize}
				\item Doesn't distinguish between error types
				\item Misleading for imbalanced datasets
				\item Doesn't account for costs of errors
			\end{itemize}
		\end{itemize}
		
		\textcolor{myred}{\textbf{EXAM TRAP:}} Accuracy = \textcolor{myblue}{percentage of all correct predictions}!
	\end{frame}
	
	\begin{frame}{Q33: Accuracy Example - Tutorial}
		\fontsize{6.5}{7.5}\selectfont
		\textbf{QUESTION: In the dividend prediction example with TP=432, TN=279, FP=121, FN=168, what is accuracy?}
		
		\vspace{0.1cm}
		\textbf{OPTIONS:}
		\begin{enumerate}
			\item 43.2\%
			\item 71.1\%
			\item 78.1\%
			\item 28.9\%
		\end{enumerate}
		
		\vspace{0.1cm}
		\textcolor{myblue}{\textbf{ANSWER: B}}
		
		\vspace{0.1cm}
		\textbf{DETAILED TUTORIAL:}
		\begin{itemize}
			\item \textbf{Calculation:}
			\begin{itemize}
				\item Accuracy = (432 + 279) / (432 + 279 + 121 + 168)
				\item = 711 / 1000
				\item = 71.1\%
			\end{itemize}
			\item \textbf{Interpretation:}
			\begin{itemize}
				\item "71.1\% of dividend predictions were correct"
				\item Error rate = 1 - 0.711 = 28.9\%
			\end{itemize}
			\item \textbf{Context:}
			\begin{itemize}
				\item 600 firms actually paid dividends
				\item 400 firms actually did not
				\item Model correctly classified 71.1\% of all firms
			\end{itemize}
		\end{itemize}
		
		\textcolor{myred}{\textbf{EXAM TRAP:}} Accuracy = \textcolor{myblue}{(TP + TN) / Total}!
	\end{frame}
	
	\begin{frame}{Q34: Error Rate Definition - Tutorial}
		\fontsize{6.5}{7.5}\selectfont
		\textbf{QUESTION: What is the error rate and how is it related to accuracy?}
		
		\vspace{0.1cm}
		\textbf{OPTIONS:}
		\begin{enumerate}
			\item Error rate = 1 - Accuracy
			\item Error rate = Accuracy + 1
			\item Error rate = Accuracy / 2
			\item Error rate = 2 × Accuracy
		\end{enumerate}
		
		\vspace{0.1cm}
		\textcolor{myblue}{\textbf{ANSWER: A}}
		
		\vspace{0.1cm}
		\textbf{DETAILED TUTORIAL:}
		\begin{itemize}
			\item \textbf{Formula:}
			\begin{itemize}
				\item $Error Rate = 1 - Accuracy = \frac{FP + FN}{TP + TN + FP + FN}$
			\end{itemize}
			\item \textbf{Interpretation:}
			\begin{itemize}
				\item "Percentage of predictions that were incorrect"
				\item Complementary to accuracy
				\item Simple and intuitive
			\end{itemize}
			\item \textbf{Example:}
			\begin{itemize}
				\item Accuracy = 71.1\%
				\item Error Rate = 100\% - 71.1\% = 28.9\%
				\item "28.9\% of predictions were wrong"
			\end{itemize}
		\end{itemize}
		
		\textcolor{myred}{\textbf{EXAM TRAP:}} Error rate = \textcolor{myblue}{1 - Accuracy}!
	\end{frame}
	
	\begin{frame}{Q35: Accuracy Limitations - Tutorial}
		\fontsize{6.5}{7.5}\selectfont
		\textbf{QUESTION: What is a major limitation of accuracy as a performance measure?}
		
		\vspace{0.1cm}
		\textbf{OPTIONS:}
		\begin{enumerate}
			\item It's too difficult to calculate
			\item It fails to distinguish between Type I and Type II errors
			\item It always favors complex models
			\item It cannot be used for imbalanced datasets
		\end{enumerate}
		
		\vspace{0.1cm}
		\textcolor{myblue}{\textbf{ANSWER: B}}
		
		\vspace{0.1cm}
		\textbf{DETAILED TUTORIAL:}
		\begin{itemize}
			\item \textbf{Limitations of accuracy:}
			\begin{enumerate}
				\item \textbf{Doesn't distinguish error types:}
				\begin{itemize}
					\item FP and FN treated equally
					\item But costs are usually different
				\end{itemize}
				\item \textbf{Misleading with imbalanced data:}
				\begin{itemize}
					\item 98\% positive, classify all as positive
					\item Accuracy = 98\%, but model is useless
				\end{itemize}
				\item \textbf{Problematic example:}
				\begin{itemize}
					\item Fraud detection: 99\% legitimate
					\item Model predicting "legitimate" always has 99\% accuracy
					\item But catches zero fraud (Recall = 0)
				\end{itemize}
			\end{enumerate}
		\end{itemize}
		
		\textcolor{myred}{\textbf{EXAM TRAP:}} Accuracy is \textcolor{myblue}{misleading for imbalanced datasets}!
	\end{frame}
	
	\begin{frame}{Q36: Precision Definition - Tutorial}
		\fontsize{6.5}{7.5}\selectfont
		\textbf{QUESTION: What is precision in classification model evaluation?}
		
		\vspace{0.1cm}
		\textbf{OPTIONS:}
		\begin{enumerate}
			\item Of all actual positives, how many were correctly identified?
			\item Of all predicted positives, how many were actually positive?
			\item Of all predictions, how many were correct?
			\item Average of true positive and true negative rates
		\end{enumerate}
		
		\vspace{0.1cm}
		\textcolor{myblue}{\textbf{ANSWER: B}}
		
		\vspace{0.1cm}
		\textbf{DETAILED TUTORIAL:}
		\begin{itemize}
			\item \textbf{Formula:}
			\begin{itemize}
				\item $Precision = \frac{TP}{TP + FP}$
			\end{itemize}
			\item \textbf{Managerial question:}
			\begin{itemize}
				\item "Of all the alerts our model generated, how many were actually real issues?"
				\item "If we act on this prediction, how often will we be right?"
			\end{itemize}
			\item \textbf{When precision matters:}
			\begin{itemize}
				\item High cost of investigation (AML alerts)
				\item Customer friction concerns
				\item Operational resource constraints
			\end{itemize}
		\end{itemize}
		
		\textcolor{myred}{\textbf{EXAM TRAP:}} Precision = \textcolor{myblue}{TP / (TP + FP)} - "of predicted positives, how many were right?"
	\end{frame}
	
	\begin{frame}{Q37: Precision Example - Tutorial}
		\fontsize{6.5}{7.5}\selectfont
		\textbf{QUESTION: In the dividend example with TP=432, FP=121, what is precision?}
		
		\vspace{0.1cm}
		\textbf{OPTIONS:}
		\begin{enumerate}
			\item 72.0\%
			\item 78.1\%
			\item 43.2\%
			\item 71.1\%
		\end{enumerate}
		
		\vspace{0.1cm}
		\textcolor{myblue}{\textbf{ANSWER: B}}
		
		\vspace{0.1cm}
		\textbf{DETAILED TUTORIAL:}
		\begin{itemize}
			\item \textbf{Calculation:}
			\begin{itemize}
				\item Precision = 432 / (432 + 121)
				\item = 432 / 553
				\item = 78.1\%
			\end{itemize}
			\item \textbf{Interpretation:}
			\begin{itemize}
				\item "When the model predicts a firm will pay a dividend, it's right 78.1\% of the time"
				\item "Of all dividend predictions, 78.1\% were correct"
			\end{itemize}
			\item \textbf{Context:}
			\begin{itemize}
				\item 432 firms correctly predicted to pay dividends
				\item 121 firms incorrectly predicted to pay dividends
				\item Total predicted positive = 553
			\end{itemize}
		\end{itemize}
		
		\textcolor{myred}{\textbf{EXAM TRAP:}} Precision = \textcolor{myblue}{TP / (TP + FP)}!
	\end{frame}
	
	\begin{frame}{Q38: Recall Definition - Tutorial}
		\fontsize{6.5}{7.5}\selectfont
		\textbf{QUESTION: What is recall (sensitivity) in classification model evaluation?}
		
		\vspace{0.1cm}
		\textbf{OPTIONS:}
		\begin{enumerate}
			\item Of all predicted positives, how many were actually positive?
			\item Of all actual positives, how many were correctly identified?
			\item Of all predictions, how many were correct?
			\item Average of true positive and true negative rates
		\end{enumerate}
		
		\vspace{0.1cm}
		\textcolor{myblue}{\textbf{ANSWER: B}}
		
		\vspace{0.1cm}
		\textbf{DETAILED TUTORIAL:}
		\begin{itemize}
			\item \textbf{Formula:}
			\begin{itemize}
				\item $Recall = \frac{TP}{TP + FN}$
			\end{itemize}
			\item \textbf{Also known as:} Sensitivity or True Positive Rate
			\item \textbf{Managerial question:}
			\begin{itemize}
				\item "Of all real fraud events that occurred, how many did we catch?"
				\item "What percentage of actual positives did we identify?"
			\end{itemize}
			\item \textbf{When recall matters:}
			\begin{itemize}
				\item Missing a positive has severe consequences
				\item Fraud detection, default prediction
				\item Compliance monitoring
			\end{itemize}
		\end{itemize}
		
		\textcolor{myred}{\textbf{EXAM TRAP:}} Recall = \textcolor{myblue}{TP / (TP + FN)} - "of actual positives, how many did we catch?"
	\end{frame}
	
	\begin{frame}{Q39: Recall Example - Tutorial}
		\fontsize{6.5}{7.5}\selectfont
		\textbf{QUESTION: In the dividend example with TP=432, FN=168, what is recall?}
		
		\vspace{0.1cm}
		\textbf{OPTIONS:}
		\begin{enumerate}
			\item 72.0\%
			\item 78.1\%
			\item 43.2\%
			\item 71.1\%
		\end{enumerate}
		
		\vspace{0.1cm}
		\textcolor{myblue}{\textbf{ANSWER: A}}
		
		\vspace{0.1cm}
		\textbf{DETAILED TUTORIAL:}
		\begin{itemize}
			\item \textbf{Calculation:}
			\begin{itemize}
				\item Recall = 432 / (432 + 168)
				\item = 432 / 600
				\item = 72.0\%
			\end{itemize}
			\item \textbf{Interpretation:}
			\begin{itemize}
				\item "Of all firms that actually paid dividends, the model caught 72\%"
				\item "The model missed 28\% of dividend payers"
			\end{itemize}
			\item \textbf{Context:}
			\begin{itemize}
				\item 600 firms actually paid dividends
				\item 432 were correctly predicted
				\item 168 were missed (False Negatives)
			\end{itemize}
		\end{itemize}
		
		\textcolor{myred}{\textbf{EXAM TRAP:}} Recall = \textcolor{myblue}{TP / (TP + FN)}!
	\end{frame}
	
	\begin{frame}{Q40: Precision vs Recall Tradeoff - Tutorial}
		\fontsize{6.5}{7.5}\selectfont
		\textbf{QUESTION: What is the relationship between precision and recall?}
		
		\vspace{0.1cm}
		\textbf{OPTIONS:}
		\begin{enumerate}
			\item They always move in the same direction
			\item There is often a tradeoff - improving one can reduce the other
			\item They are independent of each other
			\item Precision is always higher than recall
		\end{enumerate}
		
		\vspace{0.1cm}
		\textcolor{myblue}{\textbf{ANSWER: B}}
		
		\vspace{0.1cm}
		\textbf{DETAILED TUTORIAL:}
		\begin{itemize}
			\item \textbf{The tradeoff:}
			\begin{itemize}
				\item \textbf{High threshold:} Higher precision, lower recall
				\item \textbf{Low threshold:} Lower precision, higher recall
			\end{itemize}
			\item \textbf{Example:}
			\begin{itemize}
				\item Threshold = 0.9: Only very confident predictions
				\item Precision: High (few false positives)
				\item Recall: Low (miss many actual positives)
			\end{itemize}
			\item \textbf{Balancing act:}
			\begin{itemize}
				\item Business needs determine the balance
				\item F1 score captures both
				\item Cost matrix helps choose optimal point
			\end{itemize}
		\end{itemize}
		
		\textcolor{myred}{\textbf{EXAM TRAP:}} Precision and Recall have \textcolor{myblue}{tradeoff} - cannot maximize both simultaneously!
	\end{frame}
	
	\begin{frame}{Q41: F1 Score Definition - Tutorial}
		\fontsize{6.5}{7.5}\selectfont
		\textbf{QUESTION: What is the F1 Score and why is it useful?}
		
		\vspace{0.1cm}
		\textbf{OPTIONS:}
		\begin{enumerate}
			\item Sum of precision and recall
			\item Harmonic mean of precision and recall; balances both in a single metric
			\item Product of precision and recall
			\item Average of precision and recall
		\end{enumerate}
		
		\vspace{0.1cm}
		\textcolor{myblue}{\textbf{ANSWER: B}}
		
		\vspace{0.1cm}
		\textbf{DETAILED TUTORIAL:}
		\begin{itemize}
			\item \textbf{Formula:}
			\begin{itemize}
				\item $F1 = 2 \times \frac{Precision \times Recall}{Precision + Recall}$
			\end{itemize}
			\item \textbf{Why harmonic mean:}
			\begin{itemize}
				\item Harmonic mean penalizes imbalance
				\item Low precision or low recall → Low F1
				\item Better than arithmetic mean for this purpose
			\end{itemize}
			\item \textbf{Interpretation:}
			\begin{itemize}
				\item Bounded between 0 and 1
				\item Close to 1 = high precision and high recall
				\item Low score = poor precision, poor recall, or both
			\end{itemize}
		\end{itemize}
		
		\textcolor{myred}{\textbf{EXAM TRAP:}} F1 = \textcolor{myblue}{harmonic mean of precision and recall}!
	\end{frame}
	
	\begin{frame}{Q42: F1 Score Example - Tutorial}
		\fontsize{6.5}{7.5}\selectfont
		\textbf{QUESTION: Given Precision=78.1\% and Recall=72.0\%, what is the F1 Score?}
		
		\vspace{0.1cm}
		\textbf{OPTIONS:}
		\begin{enumerate}
			\item 75.0\%
			\item 74.9\%
			\item 78.1\%
			\item 72.0\%
		\end{enumerate}
		
		\vspace{0.1cm}
		\textcolor{myblue}{\textbf{ANSWER: B}}
		
		\vspace{0.1cm}
		\textbf{DETAILED TUTORIAL:}
		\begin{itemize}
			\item \textbf{Calculation:}
			\begin{itemize}
				\item F1 = 2 × (0.781 × 0.720) / (0.781 + 0.720)
				\item = 2 × 0.5623 / 1.501
				\item = 1.1246 / 1.501
				\item = 0.7494 equal 74.9\%
			\end{itemize}
			\item \textbf{Interpretation:}
			\begin{itemize}
				\item F1 = 74.9\%
				\item Balances precision and recall
				\item Good overall performance
			\end{itemize}
		\end{itemize}
		
		\textcolor{myred}{\textbf{EXAM TRAP:}} F1 equal \textcolor{myblue}{2 × (P × R) / (P + R)}!
	\end{frame}
	
	\begin{frame}{Q43: Specificity Definition - Tutorial}
		\fontsize{6.5}{7.5}\selectfont
		\textbf{QUESTION: What is specificity (true negative rate)?}
		
		\vspace{0.1cm}
		\textbf{OPTIONS:}
		\begin{enumerate}
			\item Proportion of predicted positives that are actually positive
			\item Proportion of actual negatives correctly predicted as negative
			\item Proportion of actual positives correctly predicted as positive
			\item Proportion of predictions that are correct
		\end{enumerate}
		
		\vspace{0.1cm}
		\textcolor{myblue}{\textbf{ANSWER: B}}
		
		\vspace{0.1cm}
		\textbf{DETAILED TUTORIAL:}
		\begin{itemize}
			\item \textbf{Formula:}
			\begin{itemize}
				\item $Specificity = \frac{TN}{TN + FP}$
			\end{itemize}
			\item \textbf{Interpretation:}
			\begin{itemize}
				\item "Of all actual negatives, what percentage did we correctly identify?"
				\item True Negative Rate
			\end{itemize}
			\item \textbf{Relation to Recall:}
			\begin{itemize}
				\item Recall = TPR (True Positive Rate)
				\item Specificity = TNR (True Negative Rate)
				\item Both are important for ROC curve
			\end{itemize}
		\end{itemize}
		
		\textcolor{myred}{\textbf{EXAM TRAP:}} Specificity = \textcolor{myblue}{TN / (TN + FP)} - True Negative Rate!
	\end{frame}
	
	\begin{frame}{Q44: Imbalanced Data Problem - Tutorial}
		\fontsize{6.5}{7.5}\selectfont
		\textbf{QUESTION: Why is accuracy misleading for imbalanced datasets?}
		
		\vspace{0.1cm}
		\textbf{OPTIONS:}
		\begin{enumerate}
			\item Accuracy is too difficult to calculate
			\item A model can have high accuracy by always predicting the majority class
			\item Accuracy treats all errors equally
			\item Accuracy is not a valid metric
		\end{enumerate}
		
		\vspace{0.1cm}
		\textcolor{myblue}{\textbf{ANSWER: B}}
		
		\vspace{0.1cm}
		\textbf{DETAILED TUTORIAL:}
		\begin{itemize}
			\item \textbf{The problem:}
			\begin{itemize}
				\item Dataset: 98\% solvent, 2\% default
				\item Model predicts "solvent" for everyone
				\item Accuracy = 98\% (seems great!)
				\item But model catches zero defaults
			\end{itemize}
			\item \textbf{What metrics to use instead:}
			\begin{itemize}
				\item Precision: 0\% (no defaults predicted)
				\item Recall: 0\% (caught zero actual defaults)
				\item F1: 0\% (poor precision and recall)
			\end{itemize}
			\item \textbf{Lesson:}
			\begin{itemize}
				\item Accuracy is insufficient for imbalanced data
				\item Use Precision, Recall, F1, AUC
			\end{itemize}
		\end{itemize}
		
		\textcolor{myred}{\textbf{EXAM TRAP:}} Accuracy is \textcolor{myblue}{misleading for imbalanced data}!
	\end{frame}
	
	\begin{frame}{Q45: Confusion Matrix Example - Tutorial}
		\fontsize{6.5}{7.5}\selectfont
		\textbf{QUESTION: Given a loan default model with TP=60, TN=498, FP=98, FN=5, what are the performance metrics?}
		
		\vspace{0.1cm}
		\textbf{OPTIONS:}
		\begin{enumerate}
			\item Accuracy=84.4\%, Precision=38.0\%, Recall=92.3\%
			\item Accuracy=92.3\%, Precision=38.0\%, Recall=84.4\%
			\item Accuracy=84.4\%, Precision=92.3\%, Recall=38.0\%
			\item Accuracy=38.0\%, Precision=84.4\%, Recall=92.3\%
		\end{enumerate}
		
		\vspace{0.1cm}
		\textcolor{myblue}{\textbf{ANSWER: A}}
		
		\vspace{0.1cm}
		\textbf{DETAILED TUTORIAL:}
		\begin{itemize}
			\item \textbf{Calculations:}
			\begin{itemize}
				\item Total = 60 + 498 + 98 + 5 = 661
				\item Accuracy = (60 + 498) / 661 = 558/661 = 84.4\%
				\item Precision = 60 / (60 + 98) = 60/158 = 38.0\%
				\item Recall = 60 / (60 + 5) = 60/65 = 92.3\%
			\end{itemize}
			\item \textbf{Interpretation:}
			\begin{itemize}
				\item Model catches 92.3\% of defaults (great!)
				\item But only 38\% of default predictions are correct
				\item Many false alarms (98 false positives)
			\end{itemize}
		\end{itemize}
		
		\textcolor{myred}{\textbf{EXAM TRAP:}} Work through each formula carefully - don't mix up denominators!
	\end{frame}
	
	% ============================================================
	% SECTION 4: ROC AND AUC - QUESTIONS 46-55
	% ============================================================
	
	\begin{frame}{Q46: ROC Curve Definition - Tutorial}
		\fontsize{6.5}{7.5}\selectfont
		\textbf{QUESTION: What is a Receiver Operating Characteristic (ROC) curve?}
		
		\vspace{0.1cm}
		\textbf{OPTIONS:}
		\begin{enumerate}
			\item A plot of model parameters
			\item A graphical illustration of the tradeoff between true positive and false positive rates across different thresholds
			\item A plot of training vs test performance
			\item A plot of precision vs recall
		\end{enumerate}
		
		\vspace{0.1cm}
		\textcolor{myblue}{\textbf{ANSWER: B}}
		
		\vspace{0.1cm}
		\textbf{DETAILED TUTORIAL:}
		\begin{itemize}
			\item \textbf{ROC Curve:}
			\begin{itemize}
				\item Y-axis: True Positive Rate (Recall)
				\item X-axis: False Positive Rate (1 - Specificity)
				\item Points from varying decision threshold
			\end{itemize}
			\item \textbf{Key insight:}
			\begin{itemize}
				\item Shows tradeoff between catching positives and avoiding false alarms
				\item Each threshold gives different TPR and FPR
				\item Curve connects all possible thresholds
			\end{itemize}
		\end{itemize}
		
		\textcolor{myred}{\textbf{EXAM TRAP:}} ROC shows \textcolor{myblue}{TPR vs FPR} tradeoff across thresholds!
	\end{frame}
	
	\begin{frame}{Q47: AUC Interpretation - Tutorial}
		\fontsize{6.5}{7.5}\selectfont
		\textbf{QUESTION: What does AUC (Area Under the ROC Curve) measure?}
		
		\vspace{0.1cm}
		\textbf{OPTIONS:}
		\begin{enumerate}
			\item The accuracy of the model at a specific threshold
			\item The overall ability of the model to distinguish between positive and negative classes
			\item The precision of the model
			\item The recall of the model
		\end{enumerate}
		
		\vspace{0.1cm}
		\textcolor{myblue}{\textbf{ANSWER: B}}
		
		\vspace{0.1cm}
		\textbf{DETAILED TUTORIAL:}
		\begin{itemize}
			\item \textbf{AUC Interpretation:}
			\begin{itemize}
				\item AUC = 1.0: Perfect discrimination
				\item AUC = 0.5: No discrimination (random guessing)
				\item AUC < 0.5: Worse than random (should invert predictions)
				\item Higher AUC = better model
			\end{itemize}
			\item \textbf{Probabilistic interpretation:}
			\begin{itemize}
				\item Probability that a random positive ranks higher than a random negative
				\item Measures ranking ability
				\item Threshold-independent
			\end{itemize}
		\end{itemize}
		
		\textcolor{myred}{\textbf{EXAM TRAP:}} AUC = \textcolor{myblue}{probability that positive ranks higher than negative}!
	\end{frame}
	
	\begin{frame}{Q48: AUC Values Meaning - Tutorial}
		\fontsize{6.5}{7.5}\selectfont
		\textbf{QUESTION: What does an AUC of 0.5 indicate about a model?}
		
		\vspace{0.1cm}
		\textbf{OPTIONS:}
		\begin{enumerate}
			\item Perfect predictive ability
			\item No predictive ability (random predictions)
			\item Worse than random predictions
			\item The model is overfitting
		\end{enumerate}
		
		\vspace{0.1cm}
		\textcolor{myblue}{\textbf{ANSWER: B}}
		
		\vspace{0.1cm}
		\textbf{DETAILED TUTORIAL:}
		\begin{itemize}
			\item \textbf{AUC values and interpretation:}
			\begin{itemize}
				\item AUC = 1.0: Perfect model
				\item AUC = 0.5: Random model (no predictive ability)
				\item AUC < 0.5: Worse than random (should invert)
				\item Typical good models: 0.7-0.9
			\end{itemize}
			\item \textbf{What AUC = 0.5 means:}
			\begin{itemize}
				\item Model cannot distinguish between classes
				\item Equivalent to random guessing
				\item No useful predictive information
			\end{itemize}
		\end{itemize}
		
		\textcolor{myred}{\textbf{EXAM TRAP:}} AUC = 0.5 = \textcolor{myblue}{no predictive ability} - random guessing!
	\end{frame}
	
	\begin{frame}{Q49: AUC Perfect Model - Tutorial}
		\fontsize{6.5}{7.5}\selectfont
		\textbf{QUESTION: What does an AUC of 1.0 indicate about a model?}
		
		\vspace{0.1cm}
		\textbf{OPTIONS:}
		\begin{enumerate}
			\item No predictive ability
			\item Perfect predictive ability (all positives rank above all negatives)
			\item Worse than random
			\item The model is overfitting
		\end{enumerate}
		
		\vspace{0.1cm}
		\textcolor{myblue}{\textbf{ANSWER: B}}
		
		\vspace{0.1cm}
		\textbf{DETAILED TUTORIAL:}
		\begin{itemize}
			\item \textbf{AUC = 1.0 meaning:}
			\begin{itemize}
				\item Perfect separation between classes
				\item Every positive has higher score than every negative
				\item There exists a threshold with perfect classification
				\item Usually indicates overfitting or data leakage
			\end{itemize}
			\item \textbf{Caution:}
			\begin{itemize}
				\item AUC = 1.0 on test data is suspicious
				\item Likely indicates problem with data or methodology
				\item Real-world data rarely has AUC = 1.0
			\end{itemize}
		\end{itemize}
		
		\textcolor{myred}{\textbf{EXAM TRAP:}} AUC = 1.0 = \textcolor{myblue}{perfect separation} - be suspicious!
	\end{frame}
	
	\begin{frame}{Q50: ROC Curve Shape - Tutorial}
		\fontsize{6.5}{7.5}\selectfont
		\textbf{QUESTION: Which model would have a better ROC curve?}
		
		\vspace{0.1cm}
		\textbf{OPTIONS:}
		\begin{enumerate}
			\item Model with curve closer to the upper-left corner (higher TPR, lower FPR)
			\item Model with curve closer to the diagonal line
			\item Model with curve closer to the lower-right corner
			\item All ROC curves are equally good
		\end{enumerate}
		
		\vspace{0.1cm}
		\textcolor{myblue}{\textbf{ANSWER: A}}
		
		\vspace{0.1cm}
		\textbf{DETAILED TUTORIAL:}
		\begin{itemize}
			\item \textbf{ROC curve interpretation:}
			\begin{itemize}
				\item Upper-left corner: TPR=1, FPR=0 (perfect)
				\item Diagonal line: TPR=FPR (random)
				\item Closer to upper-left = better model
			\end{itemize}
			\item \textbf{Comparison:}
			\begin{itemize}
				\item Curve A with more area = better
				\item Curve that stays above and left = better
				\item AUC captures this visually
			\end{itemize}
		\end{itemize}
		
		\textcolor{myred}{\textbf{EXAM TRAP:}} Better ROC curve = \textcolor{myblue}{closer to upper-left corner}!
	\end{frame}
	
	\begin{frame}{Q51: Threshold Effect on ROC - Tutorial}
		\fontsize{6.5}{7.5}\selectfont
		\textbf{QUESTION: What happens to ROC points as threshold changes from 0 to 1?}
		
		\vspace{0.1cm}
		\textbf{OPTIONS:}
		\begin{enumerate}
			\item Points move from (0,0) to (1,1) along the ROC curve
			\item Points move from (1,1) to (0,0) along the ROC curve
			\item Points move randomly
			\item Points stay in one place
		\end{enumerate}
		
		\vspace{0.1cm}
		\textcolor{myblue}{\textbf{ANSWER: A}}
		
		\vspace{0.1cm}
		\textbf{DETAILED TUTORIAL:}
		\begin{itemize}
			\item \textbf{Threshold extremes:}
			\begin{itemize}
				\item \textbf{Threshold = 0:} All predicted positive
				\item TPR = 1 (all positives caught)
				\item FPR = 1 (all negatives misclassified)
				\item Point: (1, 1)
			\end{itemize}
			\begin{itemize}
				\item \textbf{Threshold = 1:} All predicted negative
				\item TPR = 0 (no positives caught)
				\item FPR = 0 (all negatives correct)
				\item Point: (0, 0)
			\end{itemize}
			\item \textbf{As threshold increases:}
			\begin{itemize}
				\item Move from (1,1) toward (0,0)
				\item TPR decreases (catch fewer positives)
				\item FPR decreases (fewer false alarms)
			\end{itemize}
		\end{itemize}
		
		\textcolor{myred}{\textbf{EXAM TRAP:}} ROC curve goes from \textcolor{myblue}{(0,0) to (1,1)} as threshold decreases!
	\end{frame}
	
	\begin{frame}{Q52: AUC for Model Comparison - Tutorial}
		\fontsize{6.5}{7.5}\selectfont
		\textbf{QUESTION: How can AUC be used to compare models in lending decisions?}
		
		\vspace{0.1cm}
		\textbf{OPTIONS:}
		\begin{enumerate}
			\item Higher AUC indicates better ability to separate good from bad loans
			\item Lower AUC indicates better performance
			\item AUC is not useful for model comparison
			\item Models with same AUC perform identically
		\end{enumerate}
		
		\vspace{0.1cm}
		\textcolor{myblue}{\textbf{ANSWER: A}}
		
		\vspace{0.1cm}
		\textbf{DETAILED TUTORIAL:}
		\begin{itemize}
			\item \textbf{AUC in lending:}
			\begin{itemize}
				\item Measures ability to rank borrowers by risk
				\item Higher AUC = better separation
				\item Better model = more accurate risk ranking
			\end{itemize}
			\item \textbf{Decision context:}
			\begin{itemize}
				\item Choose model with higher AUC on test set
				\item AUC = 0.8 vs 0.7: First has better ranking
				\item Better AUC = more confidence in decisions
			\end{itemize}
		\end{itemize}
		
		\textcolor{myred}{\textbf{EXAM TRAP:}} Higher AUC = \textcolor{myblue}{better model} for lending decisions!
	\end{frame}
	
	\begin{frame}{Q53: ROC vs Precision-Recall Curve - Tutorial}
		\fontsize{6.5}{7.5}\selectfont
		\textbf{QUESTION: When would you use a Precision-Recall curve instead of ROC?}
		
		\vspace{0.1cm}
		\textbf{OPTIONS:}
		\begin{enumerate}
			\item When the dataset is perfectly balanced
			\item When the dataset is highly imbalanced (very few positives)
			\item When you only care about accuracy
			\item When you don't need to visualize performance
		\end{enumerate}
		
		\vspace{0.1cm}
		\textcolor{myblue}{\textbf{ANSWER: B}}
		
		\vspace{0.1cm}
		\textbf{DETAILED TUTORIAL:}
		\begin{itemize}
			\item \textbf{Precision-Recall curves:}
			\begin{itemize}
				\item Better for imbalanced datasets
				\item Focus on positive class performance
				\item More sensitive to rare events
			\end{itemize}
			\item \textbf{Why not ROC:}
			\begin{itemize}
				\item ROC can be optimistic with imbalance
				\item Large number of negatives makes FPR low
				\item AUC can be high even with poor positive performance
			\end{itemize}
		\end{itemize}
		
		\textcolor{myred}{\textbf{EXAM TRAP:}} Precision-Recall for \textcolor{myblue}{imbalanced data}!
	\end{frame}
	
	\begin{frame}{Q54: Multi-class ROC - Tutorial}
		\fontsize{6.5}{7.5}\selectfont
		\textbf{QUESTION: How are classification metrics extended to multi-class problems?}
		
		\vspace{0.1cm}
		\textbf{OPTIONS:}
		\begin{enumerate}
			\item Use the same formulas, just sum all classes
			\item Compute metrics for each class vs all others (one-vs-rest)
			\item Multi-class problems cannot be evaluated
			\item Use only accuracy for multi-class
		\end{enumerate}
		
		\vspace{0.1cm}
		\textcolor{myblue}{\textbf{ANSWER: B}}
		
		\vspace{0.1cm}
		\textbf{DETAILED TUTORIAL:}
		\begin{itemize}
			\item \textbf{One-vs-rest approach:}
			\begin{itemize}
				\item For each class, treat as "positive"
				\item All other classes as "negative"
				\item Compute TP, FP, FN, TN for each class
				\item Average or aggregate metrics
			\end{itemize}
			\item \textbf{Example - Credit ratings:}
			\begin{itemize}
				\item AAA, AA, A, BBB, etc.
				\item For AAA: positive = AAA, negative = all others
				\item Compute precision/recall for AAA
				\item Repeat for each rating
			\end{itemize}
		\end{itemize}
		
		\textcolor{myred}{\textbf{EXAM TRAP:}} Multi-class uses \textcolor{myblue}{one-vs-rest} approach!
	\end{frame}
	
	\begin{frame}{Q55: AUC vs Accuracy - Tutorial}
		\fontsize{6.5}{7.5}\selectfont
		\textbf{QUESTION: What is an advantage of AUC over accuracy?}
		
		\vspace{0.1cm}
		\textbf{OPTIONS:}
		\begin{enumerate}
			\item AUC is easier to calculate
			\item AUC is threshold-independent and works well for imbalanced data
			\item AUC always gives higher values
			\item AUC requires less data
		\end{enumerate}
		
		\vspace{0.1cm}
		\textcolor{myblue}{\textbf{ANSWER: B}}
		
		\vspace{0.1cm}
		\textbf{DETAILED TUTORIAL:}
		\begin{itemize}
			\item \textbf{AUC advantages:}
			\begin{itemize}
				\item \textbf{Threshold-independent:} Evaluates all thresholds
				\item \textbf{Works with imbalance:} Not skewed by majority class
				\item \textbf{Ranking ability:} Measures how well model orders cases
			\end{itemize}
			\item \textbf{Accuracy limitations:}
			\begin{itemize}
				\item Threshold-dependent (uses one threshold)
				\item Misleading for imbalanced data
				\item Doesn't capture ranking ability
			\end{itemize}
		\end{itemize}
		
		\textcolor{myred}{\textbf{EXAM TRAP:}} AUC is \textcolor{myblue}{threshold-independent} - a key advantage!
	\end{frame}
	
	% ============================================================
	% SECTION 5: THRESHOLD SELECTION - QUESTIONS 56-65
	% ============================================================
	
	\begin{frame}{Q56: Threshold Definition - Tutorial}
		\fontsize{6.5}{7.5}\selectfont
		\textbf{QUESTION: What is a decision threshold in classification?}
		
		\vspace{0.1cm}
		\textbf{OPTIONS:}
		\begin{enumerate}
			\item The point at which training stops
			\item A cutoff value applied to predicted probabilities to make binary classifications
			\item The learning rate in gradient descent
			\item The regularization parameter
		\end{enumerate}
		
		\vspace{0.1cm}
		\textcolor{myblue}{\textbf{ANSWER: B}}
		
		\vspace{0.1cm}
		\textbf{DETAILED TUTORIAL:}
		\begin{itemize}
			\item \textbf{Definition:}
			\begin{itemize}
				\item Most classification models output probabilities
				\item Threshold converts probability to class label
				\item If PD > threshold → "default"
				\item If PD less than threshold → "no default"
			\end{itemize}
			\item \textbf{Common choices:}
			\begin{itemize}
				\item Default: 0.5 (balanced)
				\item Can be adjusted based on business needs
				\item Different thresholds = different tradeoffs
			\end{itemize}
		\end{itemize}
		
		\textcolor{myred}{\textbf{EXAM TRAP:}} Threshold converts \textcolor{myblue}{probability to class label}!
	\end{frame}
	
	\begin{frame}{Q57: Cost Matrix Definition - Tutorial}
		\fontsize{6.5}{7.5}\selectfont
		\textbf{QUESTION: What is a cost matrix in threshold selection?}
		
		\vspace{0.1cm}
		\textbf{OPTIONS:}
		\begin{enumerate}
			\item A matrix showing model parameters
			\item Explicitly defined economic consequences for each quadrant of the confusion matrix
			\item A matrix showing feature correlations
			\item A matrix showing error rates
		\end{enumerate}
		
		\vspace{0.1cm}
		\textcolor{myblue}{\textbf{ANSWER: B}}
		
		\vspace{0.1cm}
		\textbf{DETAILED TUTORIAL:}
		\begin{itemize}
			\item \textbf{Cost Matrix components:}
			\begin{itemize}
				\item \textbf{False Negative:} e.g., loss from undetected fraud
				\item \textbf{False Positive:} e.g., cost to investigate alert
				\item \textbf{True Positive:} e.g., benefit of catching fraud
				\item \textbf{True Negative:} e.g., benefit of smooth processing
			\end{itemize}
			\item \textbf{Why it matters:}
			\begin{itemize}
				\item Different costs = different optimal thresholds
				\item Balances business economics with model decisions
				\item Makes threshold selection transparent
			\end{itemize}
		\end{itemize}
		
		\textcolor{myred}{\textbf{EXAM TRAP:}} Cost matrix = \textcolor{myblue}{economic consequences} of each decision!
	\end{frame}
	
	\begin{frame}{Q58: Stakeholder Input - Tutorial}
		\fontsize{6.5}{7.5}\selectfont
		\textbf{QUESTION: Who should be involved in developing a cost matrix?}
		
		\vspace{0.1cm}
		\textbf{OPTIONS:}
		\begin{enumerate}
			\item Only the modeling team
			\item Business lines, operations, finance, risk, compliance, and legal
			\item Only senior management
			\item Only regulators
		\end{enumerate}
		
		\vspace{0.1cm}
		\textcolor{myblue}{\textbf{ANSWER: B}}
		
		\vspace{0.1cm}
		\textbf{DETAILED TUTORIAL:}
		\begin{itemize}
			\item \textbf{Key stakeholders:}
			\begin{itemize}
				\item \textbf{Business lines:} Cost of missed opportunities, customer attrition
				\item \textbf{Operations:} Cost of manual reviews
				\item \textbf{Finance/Risk:} Financial impact of missed risks
				\item \textbf{Compliance/Legal:} Regulatory breach costs
			\end{itemize}
			\item \textbf{Why multiple perspectives:}
			\begin{itemize}
				\item Different stakeholders understand different costs
				\item Complete picture of all consequences
				\item Collaborative decisions are more defensible
			\end{itemize}
		\end{itemize}
		
		\textcolor{myred}{\textbf{EXAM TRAP:}} Cost matrix needs \textcolor{myblue}{cross-functional input}!
	\end{frame}
	
	\begin{frame}{Q59: Dynamic Threshold Adjustment - Tutorial}
		\fontsize{6.5}{7.5}\selectfont
		\textbf{QUESTION: Should thresholds be static or dynamic?}
		
		\vspace{0.1cm}
		\textbf{OPTIONS:}
		\begin{enumerate}
			\item Always static (set once and never change)
			\item Should be dynamic and adjusted based on changing conditions and risk appetite
			\item Always use 0.5
			\item Thresholds cannot be changed
		\end{enumerate}
		
		\vspace{0.1cm}
		\textcolor{myblue}{\textbf{ANSWER: B}}
		
		\vspace{0.1cm}
		\textbf{DETAILED TUTORIAL:}
		\begin{itemize}
			\item \textbf{Dynamic adjustment examples:}
			\begin{itemize}
				\item \textbf{Holiday fraud spike:} Lower threshold (higher recall)
				\item \textbf{Economic downturn:} Adjust default thresholds
				\item \textbf{Benign periods:} Higher threshold (fewer false positives)
			\end{itemize}
			\item \textbf{Prerequisites:}
			\begin{itemize}
				\item Clear risk appetite statement
				\item Documented review process
				\item Transparent justification
				\item Regulatory awareness
			\end{itemize}
		\end{itemize}
		
		\textcolor{myred}{\textbf{EXAM TRAP:}} Thresholds should be \textcolor{myblue}{dynamic} - not "set and forget"!
	\end{frame}
	
	\begin{frame}{Q60: Threshold Documentation - Tutorial}
		\fontsize{6.5}{7.5}\selectfont
		\textbf{QUESTION: Why should threshold choices be documented and defensible?}
		
		\vspace{0.1cm}
		\textbf{OPTIONS:}
		\begin{enumerate}
			\item Because it's a regulatory requirement
			\item To justify decisions to senior management, board, and regulators
			\item Because documentation is standard practice
			\item All of the above
		\end{enumerate}
		
		\vspace{0.1cm}
		\textcolor{myblue}{\textbf{ANSWER: D}}
		
		\vspace{0.1cm}
		\textbf{DETAILED TUTORIAL:}
		\begin{itemize}
			\item \textbf{Documentation should show:}
			\begin{itemize}
				\item How threshold aligns with risk appetite
				\item Tradeoff between recall and false positives
				\item Costs and benefits considered
				\item Periodic review process
			\end{itemize}
			\item \textbf{Example statement:}
			\begin{itemize}
				\item "Threshold set for AML monitoring to achieve 80\% recall"
				\item "Acceptable false positive rate of X\%"
				\item "Deemed acceptable given consequences of missing AML events"
			\end{itemize}
		\end{itemize}
		
		\textcolor{myred}{\textbf{EXAM TRAP:}} Threshold choices must be \textcolor{myblue}{transparent and defensible}!
	\end{frame}
	















\begin{frame}{Q61: False Positive Cost - Tutorial}
	\fontsize{6.5}{7.5}\selectfont
	
	\textbf{QUESTION: What costs are associated with False Positives in AML monitoring?}
	
	\vspace{0.1cm}
	\textbf{OPTIONS:}
	\begin{enumerate}
		\item No costs
		\item Operational cost to investigate, customer friction, resource waste
		\item Only financial costs
		\item Only reputational costs
	\end{enumerate}
	
	\vspace{0.1cm}
	\textcolor{myblue}{\textbf{ANSWER: B}}
	
	\vspace{0.1cm}
	\textbf{DETAILED TUTORIAL:}
	\begin{itemize}
		\item \textbf{False Positive costs:}
		\begin{enumerate}
			\item \textbf{Operational:} Investigator time wasted
			\item \textbf{Customer:} Friction, declined legitimate transactions
			\item \textbf{Resource:} Compliance teams chasing red herrings
			\item \textbf{Desensitization:} Alert fatigue, missing real issues
		\end{enumerate}
		
		\item \textbf{Why they matter:}
		\begin{itemize}
			\item High FP rate = inefficient operations
			\item Can outweigh benefits of catching fraud
			\item Affects customer experience
		\end{itemize}
	\end{itemize}
	
	\textcolor{myred}{\textbf{EXAM TRAP:}} False Positives = 
	\textcolor{myblue}{operational and customer costs}!
	
\end{frame}


	
	\begin{frame}{Q62: False Negative Cost - Tutorial}
		\fontsize{6.5}{7.5}\selectfont
		\textbf{QUESTION: What costs are associated with False Negatives in default prediction?}
		
		\vspace{0.1cm}
		\textbf{OPTIONS:}
		\begin{enumerate}
			\item No costs
			\item Financial losses, regulatory capital implications, missed risk
			\item Only reputational costs
			\item Only operational costs
		\end{enumerate}
		
		\vspace{0.1cm}
		\textcolor{myblue}{\textbf{ANSWER: B}}
		
		\vspace{0.1cm}
		\textbf{DETAILED TUTORIAL:}
		\begin{itemize}
			\item \textbf{False Negative costs:}
			\begin{enumerate}
				\item \textbf{Financial:} Actual loan losses
				\item \textbf{Regulatory:} Capital implications
				\item \textbf{Risk:} Misclassification of risk
				\item \textbf{Reputational:} Regulatory sanctions
			\end{enumerate}
			\item \textbf{Why they're often the largest cost:}
			\begin{itemize}
				\item Direct financial impact
				\item Hard to recover from
				\item Regulatory consequences
			\end{itemize}
		\end{itemize}
		
		\textcolor{myred}{\textbf{EXAM TRAP:}} False Negatives = \textcolor{myblue}{financial and regulatory costs} - usually the largest!
	\end{frame}
	
	\begin{frame}{Q63: Optimal Threshold Selection - Tutorial}
		\fontsize{6.5}{7.5}\selectfont
		\textbf{QUESTION: How is the optimal threshold determined using a cost matrix?}
		
		\vspace{0.1cm}
		\textbf{OPTIONS:}
		\begin{enumerate}
			\item Always use 0.5
			\item Use the threshold that minimizes total expected cost
			\item Use the threshold that maximizes accuracy
			\item Use the threshold that maximizes AUC
		\end{enumerate}
		
		\vspace{0.1cm}
		\textcolor{myblue}{\textbf{ANSWER: B}}
		
		\vspace{0.1cm}
		\textbf{DETAILED TUTORIAL:}
		\begin{itemize}
			\item \textbf{Expected cost calculation:}
			\begin{itemize}
				\item For each threshold, compute expected cost
				\item Cost = C(TP)×P(TP) + C(FP)×P(FP) + ...
				\item Choose threshold with minimum cost
			\end{itemize}
			\item \textbf{Cost-based threshold:}
			\begin{itemize}
				\item Different from accuracy-based threshold
				\item Reflects business reality
				\item More defensible
			\end{itemize}
		\end{itemize}
		
		\textcolor{myred}{\textbf{EXAM TRAP:}} Optimal threshold = \textcolor{myblue}{minimum expected cost}!
	\end{frame}
	
	\begin{frame}{Q64: Threshold and Business Needs - Tutorial}
		\fontsize{6.5}{7.5}\selectfont
		\textbf{QUESTION: Which model would be preferred for fraud detection?}
		
		\vspace{0.1cm}
		\textbf{OPTIONS:}
		\begin{enumerate}
			\item Model with high precision but low recall
			\item Model with high recall (catches most fraud) even if precision is lower
			\item Model with equal precision and recall
			\item Model with highest accuracy
		\end{enumerate}
		
		\vspace{0.1cm}
		\textcolor{myblue}{\textbf{ANSWER: B}}
		
		\vspace{0.1cm}
		\textbf{DETAILED TUTORIAL:}
		\begin{itemize}
			\item \textbf{Fraud detection priorities:}
			\begin{itemize}
				\item Missing fraud (False Negative) is very costly
				\item Better to catch more fraud (high recall)
				\item Accept some false alarms (lower precision)
				\item Business needs drive the tradeoff
			\end{itemize}
			\item \textbf{Other contexts:}
			\begin{itemize}
				\item AML monitoring: High recall critical
				\item Credit approval: Balance needed
				\item Marketing: Precision more important
			\end{itemize}
		\end{itemize}
		
		\textcolor{myred}{\textbf{EXAM TRAP:}} Fraud detection prioritizes \textcolor{myblue}{recall over precision}!
	\end{frame}
	
	\begin{frame}{Q65: Threshold Review Process - Tutorial}
		\fontsize{6.5}{7.5}\selectfont
		\textbf{QUESTION: How often should model thresholds be reviewed?}
		
		\vspace{0.1cm}
		\textbf{OPTIONS:}
		\begin{enumerate}
			\item Never - set them once
			\item Periodically, based on changing conditions and risk environment
			\item Only when regulators ask
			\item Annually without exception
		\end{enumerate}
		
		\vspace{0.1cm}
		\textcolor{myblue}{\textbf{ANSWER: B}}
		
		\vspace{0.1cm}
		\textbf{DETAILED TUTORIAL:}
		\begin{itemize}
			\item \textbf{Review triggers:}
			\begin{itemize}
				\item Significant environmental changes
				\item Shifts in risk appetite
				\item New regulatory requirements
				\item Changes in fraud patterns
			\end{itemize}
			\item \textbf{Best practice:}
			\begin{itemize}
				\item Documented review schedule
				\item Clear triggers for review
				\item Stakeholder involvement
				\item Regulatory notification when appropriate
			\end{itemize}
		\end{itemize}
		
		\textcolor{myred}{\textbf{EXAM TRAP:}} Thresholds should be \textcolor{myblue}{periodically reviewed} - not static!
	\end{frame}
	
	% ============================================================
	% SECTION 6: MODEL COMPARISON - QUESTIONS 66-80
	% ============================================================
	
	\begin{frame}{Q66: Model Comparison Example - Tutorial}
		\fontsize{6.5}{7.5}\selectfont
		\textbf{QUESTION: In the LendingClub example, how do logistic regression and neural network compare?}
		
		\vspace{0.1cm}
		\textbf{OPTIONS:}
		\begin{enumerate}
			\item Logistic regression is clearly superior
			\item Neural network is clearly superior
			\item Different metrics give different rankings; choice depends on priorities
			\item Both models perform identically on all metrics
		\end{enumerate}
		
		\vspace{0.1cm}
		\textcolor{myblue}{\textbf{ANSWER: C}}
		
		\vspace{0.1cm}
		\textbf{DETAILED TUTORIAL:}
		\begin{itemize}
			\item \textbf{Test sample comparison (Table 8.6):}
			\begin{itemize}
				\item Accuracy: LR=71.3\%, NN=70.1\% (LR better)
				\item Precision: LR=60.0\%, NN=54.1\% (LR better)
				\item Recall: LR=28.3\%, NN=37.7\% (NN better)
			\end{itemize}
			\item \textbf{Interpretation:}
			\begin{itemize}
				\item No clear winner
				\item Logistic regression: More accurate, better precision
				\item Neural network: Catches more defaults (better recall)
				\item Choice depends on business priorities
			\end{itemize}
		\end{itemize}
		
		\textcolor{myred}{\textbf{EXAM TRAP:}} Different metrics = \textcolor{myblue}{different model rankings}! Choose based on problem.
	\end{frame}
	
	\begin{frame}{Q67: Training vs Test Performance Comparison - Tutorial}
		\fontsize{6.5}{7.5}\selectfont
		\textbf{QUESTION: In the LendingClub example, what does the drop from training to test performance indicate?}
		
		\vspace{0.1cm}
		\textbf{OPTIONS:}
		\begin{enumerate}
			\item The model is perfectly specified
			\item Some overfitting is present (performance drops on unseen data)
			\item The data is perfectly clean
			\item The model needs more complexity
		\end{enumerate}
		
		\vspace{0.1cm}
		\textcolor{myblue}{\textbf{ANSWER: B}}
		
		\vspace{0.1cm}
		\textbf{DETAILED TUTORIAL:}
		\begin{itemize}
			\item \textbf{Performance drops:}
			\begin{itemize}
				\item Training accuracy: 84.2\% → Test: 71.3\% (LR)
				\item Training accuracy: 84.2\% → Test: 70.1\% (NN)
				\item Both show decrease
			\end{itemize}
			\item \textbf{Interpretation:}
			\begin{itemize}
				\item Models overfit to training data
				\item Reasonable amount of overfitting
				\item Could improve with regularization
				\item Remove irrelevant features
			\end{itemize}
		\end{itemize}
		
		\textcolor{myred}{\textbf{EXAM TRAP:}} Performance drop = \textcolor{myblue}{overfitting} - expected but should be managed!
	\end{frame}
	
	\begin{frame}{Q68: Neural Network Overfitting Signs - Tutorial}
		\fontsize{6.5}{7.5}\selectfont
		\textbf{QUESTION: How can you check for overfitting in neural networks?}
		
		\vspace{0.1cm}
		\textbf{OPTIONS:}
		\begin{enumerate}
			\item Look at training vs test performance gap
			\item Examine weights for very strong or offsetting connections
			\item Both A and B
			\item Neural networks cannot overfit
		\end{enumerate}
		
		\vspace{0.1cm}
		\textcolor{myblue}{\textbf{ANSWER: C}}
		
		\vspace{0.1cm}
		\textbf{DETAILED TUTORIAL:}
		\begin{itemize}
			\item \textbf{Overfitting checks:}
			\begin{enumerate}
				\item \textbf{Performance:} Training >> Test
				\item \textbf{Weights:} Very large or offsetting values
				\item \textbf{Loss:} Training loss flat, validation increasing
			\end{enumerate}
			\item \textbf{What to look for:}
			\begin{itemize}
				\item Large positive and negative weights
				\item Model memorizing training data
				\item Poor generalization
			\end{itemize}
		\end{itemize}
		
		\textcolor{myred}{\textbf{EXAM TRAP:}} Check \textcolor{myblue}{both performance gap and weight patterns} for overfitting!
	\end{frame}
	
	\begin{frame}{Q69: Logistic Regression Advantages - Tutorial}
		\fontsize{6.5}{7.5}\selectfont
		\textbf{QUESTION: What are advantages of logistic regression over neural networks?}
		
		\vspace{0.1cm}
		\textbf{OPTIONS:}
		\begin{enumerate}
			\item Always higher accuracy
			\item More interpretable, easier to explain, faster to train
			\item More complex and powerful
			\item Requires less data
		\end{enumerate}
		
		\vspace{0.1cm}
		\textcolor{myblue}{\textbf{ANSWER: B}}
		
		\vspace{0.1cm}
		\textbf{DETAILED TUTORIAL:}
		\begin{itemize}
			\item \textbf{Logistic regression advantages:}
			\begin{itemize}
				\item \textbf{Interpretability:} Coefficients explain feature impact
				\item \textbf{Explainability:} Easy to communicate to stakeholders
				\item \textbf{Speed:} Faster to train and deploy
				\item \textbf{Transparency:} Regulatory acceptance
			\end{itemize}
			\item \textbf{Neural network advantages:}
			\begin{itemize}
				\item Can capture complex patterns
				\item Higher potential accuracy
				\item Better for large datasets
			\end{itemize}
		\end{itemize}
		
		\textcolor{myred}{\textbf{EXAM TRAP:}} Logistic regression = \textcolor{myblue}{more interpretable} than neural networks!
	\end{frame}
	
	\begin{frame}{Q70: Neural Network Advantages - Tutorial}
		\fontsize{6.5}{7.5}\selectfont
		\textbf{QUESTION: What are advantages of neural networks over logistic regression?}
		
		\vspace{0.1cm}
		\textbf{OPTIONS:}
		\begin{enumerate}
			\item Always more interpretable
			\item Can capture complex nonlinear patterns, better for some problems
			\item Faster to train
			\item Requires less data
		\end{enumerate}
		
		\vspace{0.1cm}
		\textcolor{myblue}{\textbf{ANSWER: B}}
		
		\vspace{0.1cm}
		\textbf{DETAILED TUTORIAL:}
		\begin{itemize}
			\item \textbf{Neural network advantages:}
			\begin{itemize}
				\item \textbf{Complex patterns:} Can model nonlinear relationships
				\item \textbf{Flexibility:} Universal approximator
				\item \textbf{Performance:} Often better for complex tasks
				\item \textbf{Feature learning:} Learns representations automatically
			\end{itemize}
			\item \textbf{Disadvantages:}
			\begin{itemize}
				\item Less interpretable
				\item More data needed
				\item More computationally expensive
				\item Risk of overfitting
			\end{itemize}
		\end{itemize}
		
		\textcolor{myred}{\textbf{EXAM TRAP:}} Neural networks = \textcolor{myblue}{better for complex patterns} but less interpretable!
	\end{frame}
	
	\begin{frame}{Q71: Model Selection Criteria - Tutorial}
		\fontsize{6.5}{7.5}\selectfont
		\textbf{QUESTION: In the LendingClub example, which model would a risk manager choose?}
		
		\vspace{0.1cm}
		\textbf{OPTIONS:}
		\begin{enumerate}
			\item Always choose logistic regression
			\item Always choose neural network
			\item Depends on priorities: recall vs precision, interpretability vs accuracy
			\item Choose the model with highest accuracy
		\end{enumerate}
		
		\vspace{0.1cm}
		\textcolor{myblue}{\textbf{ANSWER: C}}
		
		\vspace{0.1cm}
		\textbf{DETAILED TUTORIAL:}
		\begin{itemize}
			\item \textbf{Tradeoffs:}
			\begin{itemize}
				\item \textbf{If catching defaults is critical:} Neural network (higher recall)
				\item \textbf{If interpretability is required:} Logistic regression
				\item \textbf{If accuracy is priority:} Logistic regression (slightly higher)
				\item \textbf{If regulatory compliance:} Logistic regression
			\end{itemize}
			\item \textbf{Key insight:}
			\begin{itemize}
				\item No single "best" model
				\item Choice depends on business needs
				\item Must align with risk appetite
			\end{itemize}
		\end{itemize}
		
		\textcolor{myred}{\textbf{EXAM TRAP:}} Model choice = \textcolor{myblue}{depends on priorities and constraints}!
	\end{frame}
	
	\begin{frame}{Q72: Overfitting Remediation - Tutorial}
		\fontsize{6.5}{7.5}\selectfont
		\textbf{QUESTION: How can overfitting be addressed in the LendingClub example?}
		
		\vspace{0.1cm}
		\textbf{OPTIONS:}
		\begin{enumerate}
			\item Add more complexity to the model
			\item Remove irrelevant features or apply regularization
			\item Use more training data only
			\item Nothing - overfitting is unavoidable
		\end{enumerate}
		
		\vspace{0.1cm}
		\textcolor{myblue}{\textbf{ANSWER: B}}
		
		\vspace{0.1cm}
		\textbf{DETAILED TUTORIAL:}
		\begin{itemize}
			\item \textbf{Remediation strategies:}
			\begin{enumerate}
				\item \textbf{Simplify model:} Remove irrelevant features
				\item \textbf{Add regularization:} L1, L2, dropout
				\item \textbf{Get more data:} More training samples
				\item \textbf{Cross-validation:} Better model selection
				\item \textbf{Early stopping:} Stop before overfitting
			\end{enumerate}
			\item \textbf{Chapter recommendation:}
			\begin{itemize}
				\item "Remove some of the least empirically relevant features"
				\item "Apply regularization to the fitted models"
			\end{itemize}
		\end{itemize}
		
		\textcolor{myred}{\textbf{EXAM TRAP:}} Fix overfitting with \textcolor{myblue}{simplification or regularization}!
	\end{frame}
	
	\begin{frame}{Q73: Model Comparison Metrics - Tutorial}
		\fontsize{6.5}{7.5}\selectfont
		\textbf{QUESTION: Which metric should be used to compare models for default prediction?}
		
		\vspace{0.1cm}
		\textbf{OPTIONS:}
		\begin{enumerate}
			\item Only accuracy
			\item A combination of metrics based on business priorities
			\item Only precision
			\item Only recall
		\end{enumerate}
		
		\vspace{0.1cm}
		\textcolor{myblue}{\textbf{ANSWER: B}}
		
		\vspace{0.1cm}
		\textbf{DETAILED TUTORIAL:}
		\begin{itemize}
			\item \textbf{Recommended approach:}
			\begin{enumerate}
				\item Start with confusion matrix
				\item Compute Accuracy, Precision, Recall, F1
				\item Consider business costs
				\item Choose metric that aligns with priorities
			\end{enumerate}
			\item \textbf{Why multiple metrics:}
			\begin{itemize}
				\item Single metric masks important differences
				\item Different stakeholders need different views
				\item Regulatory requirements may specify
			\end{itemize}
		\end{itemize}
		
		\textcolor{myred}{\textbf{EXAM TRAP:}} Use \textcolor{myblue}{multiple metrics} - don't rely on just one!
	\end{frame}
	
	\begin{frame}{Q74: Test Sample Importance - Tutorial}
		\fontsize{6.5}{7.5}\selectfont
		\textbf{QUESTION: Why is the test sample comparison more important than training sample comparison?}
		
		\vspace{0.1cm}
		\textbf{OPTIONS:}
		\begin{enumerate}
			\item Test sample is larger
			\item Test sample shows true generalization performance
			\item Training sample is always more accurate
			\item Test sample is easier to compute
		\end{enumerate}
		
		\vspace{0.1cm}
		\textcolor{myblue}{\textbf{ANSWER: B}}
		
		\vspace{0.1cm}
		\textbf{DETAILED TUTORIAL:}
		\begin{itemize}
			\item \textbf{Training sample:}
			\begin{itemize}
				\item Shows how well model memorizes
				\item Optimistically biased
				\item Not representative of new data
			\end{itemize}
			\item \textbf{Test sample:}
			\begin{itemize}
				\item Shows how well model generalizes
				\item Unbiased performance estimate
				\item What matters in practice
			\end{itemize}
		\end{itemize}
		
		\textcolor{myred}{\textbf{EXAM TRAP:}} Test sample = \textcolor{myblue}{true generalization} - always prioritize!
	\end{frame}
	
	\begin{frame}{Q75: Model Interpretation Tradeoff - Tutorial}
		\fontsize{6.5}{7.5}\selectfont
		\textbf{QUESTION: What is the tradeoff between logistic regression and neural networks?}
		
		\vspace{0.1cm}
		\textbf{OPTIONS:}
		\begin{enumerate}
			\item Logistic regression is more complex
			\item Logistic regression is more interpretable but may have lower accuracy
			\item Neural networks are always better
			\item There is no tradeoff
		\end{enumerate}
		
		\vspace{0.1cm}
		\textcolor{myblue}{\textbf{ANSWER: B}}
		
		\vspace{0.1cm}
		\textbf{DETAILED TUTORIAL:}
		\begin{itemize}
			\item \textbf{Logistic regression:}
			\begin{itemize}
				\item \textbf{Pros:} Interpretable, transparent, regulatory friendly
				\item \textbf{Cons:} Limited to linear relationships
			\end{itemize}
			\item \textbf{Neural networks:}
			\begin{itemize}
				\item \textbf{Pros:} Can capture complex patterns
				\item \textbf{Cons:} Black box, hard to explain
			\end{itemize}
			\item \textbf{Choice depends on:}
			\begin{itemize}
				\item Regulatory requirements
				\item Business needs
				\item Stakeholder expectations
			\end{itemize}
		\end{itemize}
		
		\textcolor{myred}{\textbf{EXAM TRAP:}} Tradeoff = \textcolor{myblue}{interpretability vs performance}!
	\end{frame}
	
	\begin{frame}{Q76: LendingClub Example Recap - Tutorial}
		\fontsize{6.5}{7.5}\selectfont
		\textbf{QUESTION: What is the key takeaway from the LendingClub model comparison?}
		
		\vspace{0.1cm}
		\textbf{OPTIONS:}
		\begin{enumerate}
			\item Neural networks always outperform logistic regression
			\item Logistic regression always outperforms neural networks
			\item There is no universal best model; choice depends on evaluation criteria
			\item Both models perform identically
		\end{enumerate}
		
		\vspace{0.1cm}
		\textcolor{myblue}{\textbf{ANSWER: C}}
		
		\vspace{0.1cm}
		\textbf{DETAILED TUTORIAL:}
		\begin{itemize}
			\item \textbf{Key insights:}
			\begin{itemize}
				\item Different metrics gave different rankings
				\item Logistic regression: Better accuracy and precision
				\item Neural network: Better recall
				\item No clear winner
			\end{itemize}
			\item \textbf{Implication:}
			\begin{itemize}
				\item Model choice must align with business objectives
				\item "Contradictory indicators illustrate importance of choosing appropriate evaluation metric"
				\item Regulatory context matters
			\end{itemize}
		\end{itemize}
		
		\textcolor{myred}{\textbf{EXAM TRAP:}} No universal best model - \textcolor{myblue}{choice depends on criteria}!
	\end{frame}
	
	\begin{frame}{Q77: Business Metric Alignment - Tutorial}
		\fontsize{6.5}{7.5}\selectfont
		\textbf{QUESTION: Why should evaluation metrics align with business objectives?}
		
		\vspace{0.1cm}
		\textbf{OPTIONS:}
		\begin{enumerate}
			\item Because it's a regulatory requirement
			\item To ensure model decisions reflect business priorities and risk appetite
			\item Because it makes calculations easier
			\item Because all metrics are equally important
		\end{enumerate}
		
		\vspace{0.1cm}
		\textcolor{myblue}{\textbf{ANSWER: B}}
		
		\vspace{0.1cm}
		\textbf{DETAILED TUTORIAL:}
		\begin{itemize}
			\item \textbf{Business alignment importance:}
			\begin{itemize}
				\item Metrics drive model selection
				\item Model selection drives decisions
				\item Decisions have real-world consequences
				\item Wrong metric = wrong decisions
			\end{itemize}
			\item \textbf{Example:}
			\begin{itemize}
				\item If catching fraud is priority → optimize recall
				\item If minimizing investigations → optimize precision
				\item Must match business needs
			\end{itemize}
		\end{itemize}
		
		\textcolor{myred}{\textbf{EXAM TRAP:}} Metrics must \textcolor{myblue}{align with business objectives}!
	\end{frame}
	
	\begin{frame}{Q78: Regulatory Considerations - Tutorial}
		\fontsize{6.5}{7.5}\selectfont
		\textbf{QUESTION: How do regulatory requirements affect model evaluation?}
		
		\vspace{0.1cm}
		\textbf{OPTIONS:}
		\begin{enumerate}
			\item Regulators don't care about model evaluation
			\item Regulators may specify required metrics and documentation (ECB, regulators)
			\item Only internal metrics matter
			\item Regulations only apply to training
		\end{enumerate}
		
		\vspace{0.1cm}
		\textcolor{myblue}{\textbf{ANSWER: B}}
		
		\vspace{0.1cm}
		\textbf{DETAILED TUTORIAL:}
		\begin{itemize}
			\item \textbf{Regulatory expectations:}
			\begin{itemize}
				\item ECB recommends tail-risk metrics
				\item RMSE alongside MAE
				\item Documentation of metric choice
				\item Justification for model selection
			\end{itemize}
			\item \textbf{Why it matters:}
			\begin{itemize}
				\item Compliance is mandatory
				\item Penalties for non-compliance
				\item Reputational impact
			\end{itemize}
		\end{itemize}
		
		\textcolor{myred}{\textbf{EXAM TRAP:}} Regulators \textcolor{myblue}{specify metric requirements}! 
	\end{frame}


\begin{frame}{Q79: Multiple Model Comparison - Tutorial}
	\fontsize{6.5}{7.5}\selectfont
	
	\textbf{QUESTION: How should multiple models be compared in practice?}
	
	\vspace{0.1cm}
	\textbf{OPTIONS:}
	\begin{enumerate}
		\item Choose the one with highest accuracy only
		\item Use multiple metrics, consider business context, check regulatory requirements
		\item Always choose the simplest model
		\item Always choose the most complex model
	\end{enumerate}
	
	\vspace{0.1cm}
	\textcolor{myblue}{\textbf{ANSWER: B}}
	
	\vspace{0.1cm}
	\textbf{DETAILED TUTORIAL:}
	\begin{itemize}
		\item \textbf{Comparison framework:}
		\begin{enumerate}
			\item \textbf{Multiple metrics:} Accuracy, Precision, Recall, F1, AUC
			\item \textbf{Business context:} Costs, priorities, risk appetite
			\item \textbf{Regulatory:} Required metrics, documentation
			\item \textbf{Interpretability:} Stakeholder needs
			\item \textbf{Performance:} Training vs.\ test gap
		\end{enumerate}
	\end{itemize}
	
\end{frame}


\begin{frame}{Q80: Evaluation Framework Summary - Tutorial}
	\fontsize{6.5}{7.5}\selectfont
	\textbf{QUESTION: What is the recommended framework for model performance evaluation?}
	
	\vspace{0.1cm}
	\textbf{OPTIONS:}
	\begin{enumerate}
		\item Use one metric for all problems
		\item Use multiple metrics, understand business context, document choices
		\item Always use the latest metric
		\item Only use accuracy
	\end{enumerate}
	
	\vspace{0.1cm}
	\textcolor{myblue}{\textbf{ANSWER: B}}
	
	\vspace{0.1cm}
	\textbf{DETAILED TUTORIAL:}
	\begin{itemize}
		\item \textbf{Recommended framework:}
		\begin{enumerate}
			\item \textbf{Continuous:} MSE, RMSE, MAE, MAPE
			\item \textbf{Classification:} Confusion Matrix, Accuracy, Precision, Recall, F1, AUC
			\item \textbf{Business context:} Cost matrix, risk appetite
			\item \textbf{Documentation:} Transparent and defensible
		\end{enumerate}
		\item \textbf{Key principle:}
		\begin{itemize}
			\item No single metric tells the whole story
			\item Choose based on problem and stakeholders
			\item Document and justify choices
		\end{itemize}
	\end{itemize}
	
	\textcolor{myred}{\textbf{EXAM TRAP:}} Comprehensive evaluation = \textcolor{myblue}{multiple metrics + business context + documentation}!
\end{frame}

% ============================================================
% SECTION 7: ADVANCED TOPICS - QUESTIONS 81-100
% ============================================================

\begin{frame}{Q81: Error Distribution Analysis - Tutorial}
	\fontsize{6.5}{7.5}\selectfont
	\textbf{QUESTION: What can error distributions reveal about model performance?}
	
	\vspace{0.1cm}
	\textbf{OPTIONS:}
	\begin{enumerate}
		\item Nothing beyond average metrics
		\item Systematic biases, segments with larger errors, distribution shape
		\item Only the average error
		\item Only the maximum error
	\end{enumerate}
	
	\vspace{0.1cm}
	\textcolor{myblue}{\textbf{ANSWER: B}}
	
	\vspace{0.1cm}
	\textbf{DETAILED TUTORIAL:}
	\begin{itemize}
		\item \textbf{Key insights from distributions:}
		\begin{itemize}
			\item \textbf{Bias:} Systematic over/under prediction
			\item \textbf{Segments:} Areas where model performs poorly
			\item \textbf{Shape:} Many small vs few large errors
		\end{itemize}
		\item \textbf{Visualizations to use:}
		\begin{itemize}
			\item Histograms of residuals
			\item Box plots by segment
			\item Q-Q plots for normality
		\end{itemize}
	\end{itemize}
	
	\textcolor{myred}{\textbf{EXAM TRAP:}} Error distributions reveal \textcolor{myblue}{biases and patterns} that averages hide!
\end{frame}

\begin{frame}{Q82: Systematic Bias Detection - Tutorial}
	\fontsize{6.5}{7.5}\selectfont
	\textbf{QUESTION: What does systematic bias in model errors indicate?}
	
	\vspace{0.1cm}
	\textbf{OPTIONS:}
	\begin{enumerate}
		\item The model is perfect
		\item The model consistently over-predicts or under-predicts in certain conditions
		\item The data is clean
		\item The model is overfitting
	\end{enumerate}
	
	\vspace{0.1cm}
	\textcolor{myblue}{\textbf{ANSWER: B}}
	
	\vspace{0.1cm}
	\textbf{DETAILED TUTORIAL:}
	\begin{itemize}
		\item \textbf{Systematic bias examples:}
		\begin{itemize}
			\item \textbf{Over-prediction:} Consistently predicting higher values
			\item \textbf{Under-prediction:} Consistently predicting lower values
			\item \textbf{Segment bias:} Errors in specific subgroups
		\end{itemize}
		\item \textbf{Why it matters:}
		\begin{itemize}
			\item Can lead to systematic business impact
			\item Harder to correct than random errors
			\item May indicate model misspecification
		\end{itemize}
	\end{itemize}
	
	\textcolor{myred}{\textbf{EXAM TRAP:}} Systematic bias = \textcolor{myblue}{consistent error pattern} - more problematic than random errors!
\end{frame}

\begin{frame}{Q83: Segment Analysis - Tutorial}
	\fontsize{6.5}{7.5}\selectfont
	\textbf{QUESTION: Why should performance be analyzed across different segments?}
	
	\vspace{0.1cm}
	\textbf{OPTIONS:}
	\begin{enumerate}
		\item It's not necessary
		\item To identify segments where errors are consistently larger or more biased
		\item Because all segments perform identically
		\item Because it's a regulatory requirement only
	\end{enumerate}
	
	\vspace{0.1cm}
	\textcolor{myblue}{\textbf{ANSWER: B}}
	
	\vspace{0.1cm}
	\textbf{DETAILED TUTORIAL:}
	\begin{itemize}
		\item \textbf{Segment analysis benefits:}
		\begin{itemize}
			\item Identifies model weaknesses
			\item Reveals fairness issues
			\item Highlights data quality problems
			\item Guides model improvements
		\end{itemize}
		\item \textbf{Example segments:}
		\begin{itemize}
			\item Geographic regions
			\item Product types
			\item Customer segments
			\item Economic conditions
		\end{itemize}
	\end{itemize}
	
	\textcolor{myred}{\textbf{EXAM TRAP:}} Segment analysis reveals \textcolor{myblue}{where models fail}!
\end{frame}

\begin{frame}{Q84: Risk Metrics Regulatory Preference - Tutorial}
	\fontsize{6.5}{7.5}\selectfont
	\textbf{QUESTION: Why do regulators prefer metrics that capture tail risk?}
	
	\vspace{0.1cm}
	\textbf{OPTIONS:}
	\begin{enumerate}
		\item Because tail events are rare and don't matter
		\item Because extreme events have severe financial consequences
		\item Because all errors are equally important
		\item Because they don't trust MAE
	\end{enumerate}
	
	\vspace{0.1cm}
	\textcolor{myblue}{\textbf{ANSWER: B}}
	
	\vspace{0.1cm}
	\textbf{DETAILED TUTORIAL:}
	\begin{itemize}
		\item \textbf{Tail risk importance:}
		\begin{itemize}
			\item Extreme events cause largest losses
			\item Capital requirements driven by tail risk
			\item Regulatory focus on solvency
			\item Systemic risk concerns
		\end{itemize}
		\item \textbf{Why MSE/RMSE:}
		\begin{itemize}
			\item Heavily penalize large errors
			\item Capture tail events better than MAE
			\item Align with regulatory capital frameworks
		\end{itemize}
	\end{itemize}
	
	\textcolor{myred}{\textbf{EXAM TRAP:}} Regulators want \textcolor{myblue}{tail-risk metrics} for capital adequacy!
\end{frame}

\begin{frame}{Q85: Cost Matrix Components - Tutorial}
	\fontsize{6.5}{7.5}\selectfont
	\textbf{QUESTION: What should be included in a cost matrix?}
	
	\vspace{0.1cm}
	\textbf{OPTIONS:}
	\begin{enumerate}
		\item Only financial costs
		\item All costs: financial, operational, reputational, regulatory
		\item Only operational costs
		\item Only regulatory costs
	\end{enumerate}
	
	\vspace{0.1cm}
	\textcolor{myblue}{\textbf{ANSWER: B}}
	
	\vspace{0.1cm}
	\textbf{DETAILED TUTORIAL:}
	\begin{itemize}
		\item \textbf{Cost components:}
		\begin{enumerate}
			\item \textbf{Financial:} Losses from missed risks, capital impact
			\item \textbf{Operational:} Investigation costs, resource waste
			\item \textbf{Reputational:} Customer trust, market confidence
			\item \textbf{Regulatory:} Fines, sanctions, scrutiny
		\end{enumerate}
		\item \textbf{Benefit components:}
		\begin{itemize}
			\item \textbf{True Positive:} Losses avoided, risks managed
			\item \textbf{True Negative:} Efficient processing, satisfied customers
		\end{itemize}
	\end{itemize}
	
	\textcolor{myred}{\textbf{EXAM TRAP:}} Cost matrix should capture \textcolor{myblue}{all consequences} - financial and non-financial!
\end{frame}

\begin{frame}{Q86: Threshold Justification - Tutorial}
	\fontsize{6.5}{7.5}\selectfont
	\textbf{QUESTION: How should threshold choices be justified to regulators?}
	
	\vspace{0.1cm}
	\textbf{OPTIONS:}
	\begin{enumerate}
		\item "We used the default 0.5"
		\item Through documented analysis showing alignment with risk appetite and cost-benefit tradeoffs
		\item "The model developer chose it"
		\item "We use the same threshold for all models"
	\end{enumerate}
	
	\vspace{0.1cm}
	\textcolor{myblue}{\textbf{ANSWER: B}}
	
	\vspace{0.1cm}
	\textbf{DETAILED TUTORIAL:}
	\begin{itemize}
		\item \textbf{Justification should include:}
		\begin{enumerate}
			\item Analysis of tradeoffs
			\item Cost-benefit calculations
			\item Alignment with risk appetite
			\item Stakeholder consultation
			\item Periodic review process
		\end{enumerate}
		\item \textbf{Example statement:}
		\begin{itemize}
			\item "Threshold chosen to achieve 80\% recall given severe consequences of missing fraud events"
			\item "Accepts X\% false positives costing Y in operational overhead"
		\end{itemize}
	\end{itemize}
	
	\textcolor{myred}{\textbf{EXAM TRAP:}} Thresholds must be \textcolor{myblue}{justified with documented analysis}!
\end{frame}

\begin{frame}{Q87: Stakeholder Communication - Tutorial}
	\fontsize{6.5}{7.5}\selectfont
	\textbf{QUESTION: How should model performance be communicated to different stakeholders?}
	
	\vspace{0.1cm}
	\textbf{OPTIONS:}
	\begin{enumerate}
		\item Use the same metrics for everyone
		\item Use appropriate metrics for each audience: MAPE for boards, RMSE for regulators, detailed metrics for model developers
		\item Only communicate to senior management
		\item Don't communicate model performance
	\end{enumerate}
	
	\vspace{0.1cm}
	\textcolor{myblue}{\textbf{ANSWER: B}}
	
	\vspace{0.1cm}
	\textbf{DETAILED TUTORIAL:}
	\begin{itemize}
		\item \textbf{Audience considerations:}
		\begin{itemize}
			\item \textbf{Board/Management:} MAPE, Accuracy (percentage-based)
			\item \textbf{Regulators:} RMSE with MAE, tail-risk metrics
			\item \textbf{Model Developers:} All metrics, error distributions
			\item \textbf{Audit Committee:} MAPE for standardized comparison
		\end{itemize}
		\item \textbf{Communication best practices:}
		\begin{itemize}
			\item Use visualizations (ROC curves, error distributions)
			\item Explain what metrics mean in business terms
			\item Highlight tradeoffs and limitations
		\end{itemize}
	\end{itemize}
	
	\textcolor{myred}{\textbf{EXAM TRAP:}} Different stakeholders need \textcolor{myblue}{different metrics and explanations}!
\end{frame}

\begin{frame}{Q88: Performance Monitoring Frequency - Tutorial}
	\fontsize{6.5}{7.5}\selectfont
	\textbf{QUESTION: How often should model performance be monitored?}
	
	\vspace{0.1cm}
	\textbf{OPTIONS:}
	\begin{enumerate}
		\item Only during model development
		\item Continuously or at regular intervals (e.g., quarterly, annually)
		\item Only when regulators ask
		\item Never after deployment
	\end{enumerate}
	
	\vspace{0.1cm}
	\textcolor{myblue}{\textbf{ANSWER: B}}
	
	\vspace{0.1cm}
	\textbf{DETAILED TUTORIAL:}
	\begin{itemize}
		\item \textbf{Monitoring frequency:}
		\begin{itemize}
			\item \textbf{Regular intervals:} Quarterly, semi-annually, annually
			\item \textbf{Event-driven:} Significant environmental changes
			\item \textbf{Trigger-based:} Performance threshold breaches
		\end{itemize}
		\item \textbf{What to monitor:}
		\begin{itemize}
			\item Performance metrics over time
			\item Error distributions
			\item Segment performance
			\item Threshold appropriateness
		\end{itemize}
	\end{itemize}
	
	\textcolor{myred}{\textbf{EXAM TRAP:}} Model performance must be \textcolor{myblue}{monitored continuously} - not just during development!
\end{frame}

\begin{frame}{Q89: Model Validation Process - Tutorial}
	\fontsize{6.5}{7.5}\selectfont
	\textbf{QUESTION: What is the role of validation in model performance evaluation?}
	
	\vspace{0.1cm}
	\textbf{OPTIONS:}
	\begin{enumerate}
		\item Validation is the same as training
		\item Validation provides independent assessment of model performance
		\item Validation is only for data preparation
		\item Validation is not necessary
	\end{enumerate}
	
	\vspace{0.1cm}
	\textcolor{myblue}{\textbf{ANSWER: B}}
	
	\vspace{0.1cm}
	\textbf{DETAILED TUTORIAL:}
	\begin{itemize}
		\item \textbf{Validation purposes:}
		\begin{enumerate}
			\item Independent performance assessment
			\item Hyperparameter tuning
			\item Overfitting detection
			\item Model comparison
		\end{enumerate}
		\item \textbf{Validation best practices:}
		\begin{itemize}
			\item Use separate validation data
			\item Use cross-validation
			\item Document validation process
			\item Regulatory compliance
		\end{itemize}
	\end{itemize}
	
	\textcolor{myred}{\textbf{EXAM TRAP:}} Validation provides \textcolor{myblue}{independent assessment} of model performance!
\end{frame}

\begin{frame}{Q90: Performance Metrics Documentation - Tutorial}
	\fontsize{6.5}{7.5}\selectfont
	\textbf{QUESTION: Why should performance metric choices be documented?}
	
	\vspace{0.1cm}
	\textbf{OPTIONS:}
	\begin{enumerate}
		\item It's not necessary
		\item For transparency, regulatory compliance, and model governance
		\item Only if regulators ask
		\item Because it's standard practice
	\end{enumerate}
	
	\vspace{0.1cm}
	\textcolor{myblue}{\textbf{ANSWER: B}}
	
	\vspace{0.1cm}
	\textbf{DETAILED TUTORIAL:}
	\begin{itemize}
		\item \textbf{Documentation should include:}
		\begin{enumerate}
			\item Choice of metrics
			\item Justification for choices
			\item Business context
			\item Regulatory requirements
			\item Performance results
			\item Model selection rationales
		\end{enumerate}
		\item \textbf{Why it matters:}
		\begin{itemize}
			\item Regulatory compliance
			\item Audit trail
			\item Knowledge transfer
			\item Defensible decisions
		\end{itemize}
	\end{itemize}
	
	\textcolor{myred}{\textbf{EXAM TRAP:}} Metric choices must be \textcolor{myblue}{documented and justified}!
\end{frame}

\begin{frame}{Q91: Performance Decay Detection - Tutorial}
	\fontsize{6.5}{7.5}\selectfont
	\textbf{QUESTION: How can model performance decay be detected?}
	
	\vspace{0.1cm}
	\textbf{OPTIONS:}
	\begin{enumerate}
		\item It cannot be detected
		\item By monitoring performance metrics over time and comparing to baseline
		\item Only when models fail catastrophically
		\item By monitoring training performance only
	\end{enumerate}
	
	\vspace{0.1cm}
	\textcolor{myblue}{\textbf{ANSWER: B}}
	
	\vspace{0.1cm}
	\textbf{DETAILED TUTORIAL:}
	\begin{itemize}
		\item \textbf{Detection methods:}
		\begin{itemize}
			\item Track metrics over time
			\item Compare to baseline performance
			\item Monitor error distributions
			\item Segment performance tracking
		\end{itemize}
		\item \textbf{Common causes:}
		\begin{itemize}
			\item Data drift (concept drift)
			\item Economic changes
			\item Regulatory changes
			\item Model degradation
		\end{itemize}
	\end{itemize}
	
	\textcolor{myred}{\textbf{EXAM TRAP:}} Monitor metrics over time to detect \textcolor{myblue}{performance decay}!
\end{frame}

\begin{frame}{Q92: Concept Drift - Tutorial}
	\fontsize{6.5}{7.5}\selectfont
	\textbf{QUESTION: What is concept drift and how does it affect model performance?}
	
	\vspace{0.1cm}
	\textbf{OPTIONS:}
	\begin{enumerate}
		\item It doesn't affect model performance
		\item Changes in the underlying relationship between inputs and outputs over time
		\item Only affects training data
		\item Only affects test data
	\end{enumerate}
	
	\vspace{0.1cm}
	\textcolor{myblue}{\textbf{ANSWER: B}}
	
	\vspace{0.1cm}
	\textbf{DETAILED TUTORIAL:}
	\begin{itemize}
		\item \textbf{Concept drift examples:}
		\begin{itemize}
			\item Changing fraud patterns
			\item Economic cycles affecting default relationships
			\item Customer behavior changes
			\item Regulatory changes
		\end{itemize}
		\item \textbf{Detection and response:}
		\begin{itemize}
			\item Monitor performance metrics
			\item Regularly retrain models
			\item Update model parameters
			\item Review thresholds
		\end{itemize}
	\end{itemize}
	
	\textcolor{myred}{\textbf{EXAM TRAP:}} Concept drift = \textcolor{myblue}{changing relationships} over time - requires monitoring!
\end{frame}

\begin{frame}{Q93: Model Retraining Triggers - Tutorial}
	\fontsize{6.5}{7.5}\selectfont
	\textbf{QUESTION: What should trigger model retraining?}
	
	\vspace{0.1cm}
	\textbf{OPTIONS:}
	\begin{enumerate}
		\item Never - models are permanent
		\item Performance decay below thresholds, significant data changes, or scheduled intervals
		\item Only when regulators require
		\item Randomly
	\end{enumerate}
	
	\vspace{0.1cm}
	\textcolor{myblue}{\textbf{ANSWER: B}}
	
	\vspace{0.1cm}
	\textbf{DETAILED TUTORIAL:}
	\begin{itemize}
		\item \textbf{Retraining triggers:}
		\begin{enumerate}
			\item \textbf{Performance:} Metrics below acceptable levels
			\item \textbf{Data:} Significant changes in data distribution
			\item \textbf{Schedule:} Regular intervals (monthly, quarterly)
			\item \textbf{Events:} Major economic or regulatory changes
		\end{enumerate}
		\item \textbf{Best practice:}
		\begin{itemize}
			\item Document retraining triggers
			\item Regular performance review
			\item Proactive monitoring
			\item Automated alerts
		\end{itemize}
	\end{itemize}
	
	\textcolor{myred}{\textbf{EXAM TRAP:}} Retrain when \textcolor{myblue}{performance decays, data changes, or on schedule}!
\end{frame}

\begin{frame}{Q94: Performance Reporting - Tutorial}
	\fontsize{6.5}{7.5}\selectfont
	\textbf{QUESTION: What should be included in model performance reports?}
	
	\vspace{0.1cm}
	\textbf{OPTIONS:}
	\begin{enumerate}
		\item Only the latest metrics
		\item Current metrics, trends, comparisons to baseline, and any issues identified
		\item Only when performance is poor
		\item Only for regulatory submissions
	\end{enumerate}
	
	\vspace{0.1cm}
	\textcolor{myblue}{\textbf{ANSWER: B}}
	
	\vspace{0.1cm}
	\textbf{DETAILED TUTORIAL:}
	\begin{itemize}
		\item \textbf{Report components:}
		\begin{enumerate}
			\item \textbf{Current:} Latest performance metrics
			\item \textbf{Trends:} Performance over time
			\item \textbf{Baseline:} Comparison to initial/expected performance
			\item \textbf{Issues:} Any identified problems or concerns
			\item \textbf{Actions:} Planned or taken remediation
		\end{enumerate}
		\item \textbf{Audience:}
		\begin{itemize}
			\item Model owners
			\item Management
			\item Regulators
			\item Audit committees
		\end{itemize}
	\end{itemize}
	
	\textcolor{myred}{\textbf{EXAM TRAP:}} Reports should include \textcolor{myblue}{current metrics, trends, and issues}!
\end{frame}

\begin{frame}{Q95: Fairness and Bias - Tutorial}
	\fontsize{6.5}{7.5}\selectfont
	\textbf{QUESTION: How does model performance evaluation help identify fairness issues?}
	
	\vspace{0.1cm}
	\textbf{OPTIONS:}
	\begin{enumerate}
		\item It doesn't - fairness is separate
		\item By analyzing performance across different demographic or customer segments
		\item Only if fairness metrics are explicitly included
		\item Fairness cannot be measured
	\end{enumerate}
	
	\vspace{0.1cm}
	\textcolor{myblue}{\textbf{ANSWER: B}}
	
	\vspace{0.1cm}
	\textbf{DETAILED TUTORIAL:}
	\begin{itemize}
		\item \textbf{Fairness evaluation:}
		\begin{itemize}
			\item Compare performance across segments
			\item Check for systematic bias
			\item Ensure equal treatment
			\item Regulatory compliance (fair lending)
		\end{itemize}
		\item \textbf{What to examine:}
		\begin{itemize}
			\item Accuracy differences
			\item False positive/negative rates
			\item Disparate impact
			\item Equal opportunity
		\end{itemize}
	\end{itemize}
	
	\textcolor{myred}{\textbf{EXAM TRAP:}} Segment analysis reveals \textcolor{myblue}{fairness issues}!
\end{frame}

\begin{frame}{Q96: Explainable AI (XAI) - Tutorial}
	\fontsize{6.5}{7.5}\selectfont
	\textbf{QUESTION: How does model performance evaluation relate to explainable AI?}
	
	\vspace{0.1cm}
	\textbf{OPTIONS:}
	\begin{enumerate}
		\item It doesn't - they are unrelated
		\item Performance evaluation identifies where models need better explanations
		\item XAI is only for regulators
		\item XAI replaces performance evaluation
	\end{enumerate}
	
	\vspace{0.1cm}
	\textcolor{myblue}{\textbf{ANSWER: B}}
	
	\vspace{0.1cm}
	\textbf{DETAILED TUTORIAL:}
	\begin{itemize}
		\item \textbf{Connection to XAI:}
		\begin{itemize}
			\item Poor performance in segments → need explanations
			\item Complex models (DNNs) → need XAI techniques
			\item Regulatory requirements for explanation
			\item Tools: SHAP, LIME for interpretation
		\end{itemize}
		\item \textbf{Practical application:}
		\begin{itemize}
			\item Use explanations to understand poor performance
			\item Identify feature importance
			\item Debug model issues
			\item Build trust in predictions
		\end{itemize}
	\end{itemize}
	
	\textcolor{myred}{\textbf{EXAM TRAP:}} Performance evaluation guides \textcolor{myblue}{where explanations are needed}!
\end{frame}

\begin{frame}{Q97: SHAP Values for Interpretability - Tutorial}
	\fontsize{6.5}{7.5}\selectfont
	\textbf{QUESTION: How do SHAP values help with model performance evaluation?}
	
	\vspace{0.1cm}
	\textbf{OPTIONS:}
	\begin{enumerate}
		\item They don't - they only explain predictions
		\item They help identify which features are driving good and bad performance
		\item They replace performance metrics
		\item They only work for linear models
	\end{enumerate}
	
	\vspace{0.1cm}
	\textcolor{myblue}{\textbf{ANSWER: B}}
	
	\vspace{0.1cm}
	\textbf{DETAILED TUTORIAL:}
	\begin{itemize}
		\item \textbf{SHAP in evaluation:}
		\begin{itemize}
			\item Identify important features
			\item Understand model behavior
			\item Debug performance issues
			\item Explain to stakeholders
		\end{itemize}
		\item \textbf{Application:}
		\begin{itemize}
			\item Feature importance analysis
			\item Understanding model decisions
			\item Regulatory explanations
			\item Model improvement insights
		\end{itemize}
	\end{itemize}
	
	\textcolor{myred}{\textbf{EXAM TRAP:}} SHAP helps \textcolor{myblue}{understand what drives model performance}!
\end{frame}

\begin{frame}{Q98: LIME for Local Explanations - Tutorial}
	\fontsize{6.5}{7.5}\selectfont
	\textbf{QUESTION: When is LIME particularly useful for model evaluation?}
	
	\vspace{0.1cm}
	\textbf{OPTIONS:}
	\begin{enumerate}
		\item For evaluating overall model performance
		\item For explaining individual predictions and understanding edge cases
		\item For model training
		\item For data preparation
	\end{enumerate}
	
	\vspace{0.1cm}
	\textcolor{myblue}{\textbf{ANSWER: B}}
	
	\vspace{0.1cm}
	\textbf{DETAILED TUTORIAL:}
	\begin{itemize}
		\item \textbf{LIME applications:}
		\begin{itemize}
			\item Explain specific predictions
			\item Understand why model made certain decisions
			\item Identify problematic cases
			\item Build trust with stakeholders
		\end{itemize}
		\item \textbf{Evaluation insights:}
		\begin{itemize}
			\item Why certain cases were misclassified
			\item What features drove errors
			\item Edge case behavior
			\item Model limitations
		\end{itemize}
	\end{itemize}
	
	\textcolor{myred}{\textbf{EXAM TRAP:}} LIME explains \textcolor{myblue}{individual predictions} - useful for error analysis!
\end{frame}

\begin{frame}{Q99: Model Governance - Tutorial}
	\fontsize{6.5}{7.5}\selectfont
	\textbf{QUESTION: What is the role of model governance in performance evaluation?}
	
	\vspace{0.1cm}
	\textbf{OPTIONS:}
	\begin{enumerate}
		\item No role - models govern themselves
		\item Ensures consistent evaluation, documentation, and monitoring across all models
		\item Only for regulatory compliance
		\item Only for complex models
	\end{enumerate}
	
	\vspace{0.1cm}
	\textcolor{myblue}{\textbf{ANSWER: B}}
	
	\vspace{0.1cm}
	\textbf{DETAILED TUTORIAL:}
	\begin{itemize}
		\item \textbf{Governance components:}
		\begin{enumerate}
			\item \textbf{Standardization:} Consistent metrics across models
			\item \textbf{Documentation:} Clear records of evaluation
			\item \textbf{Monitoring:} Regular performance review
			\item \textbf{Oversight:} Independent validation
		\end{enumerate}
		\item \textbf{Why it matters:}
		\begin{itemize}
			\item Regulatory compliance
			\item Risk management
			\item Consistent decision-making
			\item Auditability
		\end{itemize}
	\end{itemize}
	
	\textcolor{myred}{\textbf{EXAM TRAP:}} Governance ensures \textcolor{myblue}{consistent evaluation and monitoring}!
\end{frame}

\begin{frame}{Q100: Chapter Summary - Tutorial}
	\fontsize{6.5}{7.5}\selectfont
	\textbf{QUESTION: What are the key takeaways from model performance evaluation?}
	
	\vspace{0.1cm}
	\textbf{OPTIONS:}
	\begin{enumerate}
		\item Only accuracy matters
		\item Use multiple metrics, consider business context, monitor continuously, document thoroughly
		\item Only use what regulators require
		\item Evaluation is only needed during development
	\end{enumerate}
	
	\vspace{0.1cm}
	\textcolor{myblue}{\textbf{ANSWER: B}}
	
	\vspace{0.1cm}
	\textbf{DETAILED TUTORIAL:}
	\begin{itemize}
		\item \textbf{Key takeaways:}
		\begin{enumerate}
			\item \textbf{Multiple metrics:} Use continuous (MSE, RMSE, MAE, MAPE) and classification (Confusion Matrix, Accuracy, Precision, Recall, F1, AUC)
			\item \textbf{Business context:} Align with business objectives and risk appetite
			\item \textbf{Continuous monitoring:} Track performance over time
			\item \textbf{Documentation:} Record and justify all choices
		\end{enumerate}
		\item \textbf{Critical insights:}
		\begin{itemize}
			\item No single metric tells the whole story
			\item Different metrics serve different purposes
			\item Performance evaluation is a continuous process
			\item Regulatory requirements must be met
		\end{itemize}
	\end{itemize}
	
	\textcolor{myred}{\textbf{EXAM TRAP:}} Comprehensive evaluation = \textcolor{myblue}{multiple metrics + business context + monitoring + documentation}!
\end{frame}

% ============================================================
% SUMMARY SLIDE
% ============================================================

\begin{frame}{Summary: Complete Review}
	\fontsize{6.5}{7.5}\selectfont
	\textbf{\textcolor{myblue}{Continuous Output Metrics:}}
	\begin{itemize}
		\item \textcolor{myblue}{MSE:} Average squared error; penalizes large errors; captures tail risk
		\item \textcolor{myblue}{RMSE:} Square root of MSE; same units as target; easier to interpret
		\item \textcolor{myblue}{MAE:} Average absolute error; robust to outliers
		\item \textcolor{myblue}{MAPE:} Percentage error; board-friendly; compares across scales
	\end{itemize}
	
	\vspace{0.1cm}
	\textbf{\textcolor{myblue}{Classification Metrics:}}
	\begin{itemize}
		\item \textcolor{myblue}{Confusion Matrix:} TP, TN, FP, FN - foundation of all metrics
		\item \textcolor{myblue}{Accuracy:} Total correct / total - misleading for imbalanced data
		\item \textcolor{myblue}{Precision:} TP / (TP + FP) - "of predicted positives, how many were right?"
		\item \textcolor{myblue}{Recall:} TP / (TP + FN) - "of actual positives, how many did we catch?"
		\item \textcolor{myblue}{F1:} Harmonic mean of precision and recall
		\item \textcolor{myblue}{AUC:} Area under ROC curve - threshold-independent ranking ability
	\end{itemize}
	
	\vspace{0.1cm}
	\textbf{\textcolor{myblue}{Threshold Selection:}}
	\begin{itemize}
		\item \textcolor{myblue}{Cost Matrix:} Economic consequences of each decision
		\item \textcolor{myblue}{Tradeoff:} Balance false positives vs false negatives
		\item \textcolor{myblue}{Dynamic:} Adjust based on changing conditions
		\item \textcolor{myblue}{Documentation:} Transparent and defensible
	\end{itemize}
	
	\vspace{0.1cm}
	\textcolor{myred}{\textbf{Exam Success:}} Understand the metrics, their tradeoffs, and when to use each!
\end{frame}

\end{document}